% !TEX root = ../../tree_notes.tex
% =====================================================================
%  Chapter 25 -- Link-Cut Trees (Dynamic Trees)                   WRITTEN
%  (folder ch23_link_cut_trees; label prefix c23; PRINTS as
%  Chapter 25; folder number 23; build with:  python tools/build.py --only 23)
%
%  NUMBERING TRAP (the Ch.22/23/24 trap, mirrored): \chref{c23} prints
%  "Chapter 25" (THIS chapter).  Do NOT confuse the label with the printed
%  number: \chref{c16} prints "Chapter 22" (HLD), \chref{c29} prints
%  "Chapter 23" (Li Chao), \chref{c18} prints "Chapter 24" (centroid),
%  \chref{c30} prints "Chapter 26" (virtual trees), \chref{c21} prints
%  "Chapter 27" (Mo), \chref{c06} prints "Chapter 8" (balanced trees / the
%  splay INTRODUCTION).  Every cite of THIS chapter is \chref{c23}.
%
%  Owns DYNAMIC TREES -- the splay engine (rotations driven by access, the
%  amortised O(log n) splaying analysis, top-down splaying, subtree-size
%  aggregates), the LINK-CUT TREE (preferred paths, access, makeroot/evert with
%  the reverse lazy tag, link, cut, path aggregates), and the boundary of what a
%  dynamic tree can and cannot maintain (virtual children -> subtree aggregates,
%  Euler-tour trees for connectivity vs LCT for path queries).
%  It STARTS where Ch.8 8.5 ENDS: \secref{c06-splay} owns the splay DEFINITION
%  and the amortised EXCHANGE-ARGUMENT-lite intro and EXPLICITLY promises "full
%  analysis, top-down splaying, and link-cut dynamic trees are Chapter 25."
%  THIS chapter discharges that promise.  It never re-derives the BST rotation
%  count or the binary-search-tree facts -- those live in Ch.5/Ch.6/Ch.8.
%  Sections:
%    * 25.1 the splay engine -- splaying by access, zig / zig-zig / zig-zag,
%      the potential method giving amortised O(log n) per operation, top-down
%      splaying (no recursion, no parent pointers needed), and size-maintained
%      splay trees as an order-statistics / implicit-index sequence structure.
%      DISCHARGES the \secref{c06-splay} forward-cite.  Solved SPOJ ORDERSET
%      (order-statistics set via a size-maintained splay) -- insert/delete/
%      k-th/count-less-than on a dynamic set, |x| <= 1e9, Q <= 2e5.
%    * 25.2 preferred paths and access -- a rooted forest cut into heavy
%      (preferred) paths, each stored as one aux SPLAY tree keyed by depth;
%      the access(v) spine walk; the path-parent pointer; is_root; and reading
%      the LCA of two nodes with two accesses (no makeroot).  Solved SPOJ
%      DYNALCA (dynamic forest, link A B / cut A / lca A B, N,M <= 1e5) -- the
%      canonical rooted-LCT exercise, hand-replayed to 1,2,3,2.
%    * 25.3 evert, link, cut, path aggregates -- makeroot(v) = access + reverse
%      lazy tag so any node becomes the root; link/cut; and pushing a monoid
%      (sum / min / max over VERTEX values) up the aux trees to answer a path
%      query in O(log n).  WORKED example: dynamic path vertex-sum (the
%      smallest honest LCT you can write), fully traced.  NO solvedproblem here
%      (the in-section solved count stays 2; the flagship lives at the Summit).
%    * 25.4 beyond paths -- the hard part of a real LCT: a node's SPLAY children
%      are only its PREFERRED heavy children, so a SUBTREE aggregate needs a
%      separate virtual-children accumulator (vsum) that is fixed up whenever an
%      access detaches a solid child; the connectivity question and why an
%      EULER-TOUR TREE (or DSU / rollback DSU) is often the simpler tool; the
%      dynamic spanning forest idea.  Capstone WORKED examples + the LCT-vs-ETT
%      pick table def:c23-map.  NO solvedproblem.
%    * 25.5 the engine boundary -- when a link-cut tree is the WRONG tool:
%      static path RANGE queries on a FIXED tree -> HLD (Ch.22 c16, simpler,
%      faster constant); subtree queries on a fixed tree -> Euler + Fenwick
%      (Ch.5 c24 / Ch.10); fully offline edge add/rollback -> DSU with rollback
%      or divide-and-conquer (Ch.24 engine); only connectivity -> ETT/DSU.
%      WORKED examples only -- NO solvedproblem (in-section count stays 2).
%  Chapter-end Summit: CF 487E "Tourists" (3200, L6) -- a general GRAPH with a
%  price on each vertex, operations C a w (change price) and A a b (over ALL
%  simple paths a->b, the minimum vertex price you are forced to visit).
%  Condense the biconnected components into the BLOCK-CUT TREE, put each
%  component's price multiset minimum on its component node, then an LCT path-
%  min over the block-cut tree answers A and updates on C.  Both samples fully
%  hand-replayed (S1 triangle -> 1,2; S2 seven-node chain of triangles -> 2,
%  1,5,3), including why A 2 3 = 2 in S2 (the BCT path avoids component {1,2,5}).
%
%  Diagram-first: this chapter is TREE based.  The running example is the shared
%  rooted forest R used across sections (defined once, authoritatively, in
%  00_opener).  Rooted at 1, edges (parent -> child, 1-based):
%    1->2, 1->3, 2->4, 2->5, 5->6, 5->7, 3->8
%  (8 nodes; a star-of-stars shape).  The Summit uses its OWN tiny graphs (the
%  CF samples), never R, so no trace collides.
%
%  Anti-repetition contract (teach once, cross-reference by label):
%    * Ch.8 8.5 (\chref{c06}) owns the splay DEFINITION and the intro amortised
%      argument; 25.1 REUSES that vocabulary and only adds the full potential
%      analysis + top-down splaying + size aggregates -- the single rotation is
%      never re-drawn.
%    * Ch.22 HLD (\chref{c16}) owns static heavy-light path decomposition; the
%      LCT "preferred path" is the SAME idea made dynamic -- 25.2 says so once
%      and back-cites \secref{c16-*}; 25.5 draws the "static -> use HLD" line.
%    * Ch.5 Euler tour (\chref{c24}) owns the flatten-to-array trick; 25.4 cites
%      it for the ETT comparison instead of re-teaching tin/tout.
%    * Ch.16 LCA (\chref{c11}) owns \lca and the binary-lifting tables; 25.2's
%      access-LCA is an ALTERNATIVE, not a re-derivation, and back-cites it.
%    * Ch.24 centroid (\chref{c18}) owns the offline D&C-on-trees engine; 25.5
%      routes "offline edge timeline" problems to it (and to rollback DSU), never
%      re-solving them with a full LCT.
%    * Explicitly OUT (prose + forward-cite only): static path ranges -> HLD
%      Ch.22 (\chref{c16}); fixed-tree subtree ranges -> Euler/Fenwick Ch.5/10;
%      offline dynamic connectivity -> rollback DSU; a subset of marked nodes
%      compressed to its Steiner tree -> virtual trees Ch.26 (\chref{c30}).
%
%  Known URL collisions (confirmed by grep of chapters/ + the live CF API;
%  prose-cite only, NEVER re-link -- urlsweep must print an empty DUPLICATES
%  section):
%    * CF 343D "Water Tree" (2100) -> Ch.8 (\chref{c06}) 8.5 practice.  The
%      classic splay/HLD plumbing exercise; its URL lives in Ch.8, PROSE only.
%    * CF 1416D "Graph and Queries" (2900) -> Ch.19 (\chref{c19}) persistent
%      structures.  A max-flow-forest + segment-tree-of-sets task, NOT an LCT;
%      never re-linked here.
%    * CF 1601C "Optimal Insertion" -> Ch.19 practice. Untouched by c23.
%    * The whole Ch.24 centroid roster (321C/715C/293E/833D/757G) and the Ch.22
%      HLD roster are untouched.  SPOJ ORDERSET and SPOJ DYNALCA are the only
%      solved-problem URLs and appear in NO other chapter.
%
%  Roster (every title/rating/tag/sample live-verified against the Codeforces
%  problemset API / problem page / SPOJ BEFORE a trace was written; CF-forward):
%    solved: SPOJ ORDERSET  OrderStatisticTree     (L4) (25.1) -- replayed 1,2,2,invalid
%            SPOJ DYNALCA   Dynamic LCA            (L5) (25.2) -- replayed 1,2,3,2
%    worked: a size-maintained splay insert/split (25.1); the access spine on R
%            (25.2); makeroot + reverse-tag on a 4-node path (25.3); the dynamic
%            path vertex-sum trace (25.3); the virtual-children vsum fix-up
%            (25.4); the LCT-vs-ETT connectivity pick (25.4); the engine pick
%            table (25.5).
%    summit: CF 487E Tourists                     (3200, L6) -- general graph,
%            vertex prices, C a w change price / A a b = minimum price forced on
%            every simple a->b path; block-cut tree condensation + LCT path-min
%            with a multiset per biconnected component.  Both samples
%            hand-replayed (S1 -> 1,2; S2 -> 2,1,5,3).
%    mixed (10, ascending by live-verified rating, every ID confirmed ABSENT from
%            every other chapter via the grep sweep -- locked at the
%            99_chapter_end write): 1975E(2100), 1408E(2400), 1738F(2400),
%            696E(3000), 603E(3000), 1628E(3100), 809E(3100), 1019E(3200),
%            1254E(3300), 1548E(3400).  Most are DELIBERATE REJECTION drills --
%            their hints name the TRUE tool (rollback DSU, HLD, centroid, Euler)
%            that beats an LCT.
% =====================================================================
\bookchapter{Link-Cut Trees (Dynamic Trees)}{c23}
\chapterinfo{CF 2300--2800}{Written}

% =====================================================================
%  Chapter 25 opener -- "Chapter at a Glance", three reading paths + the
%  ladder staircase, THE running example (the shared 8-node rooted forest R),
%  the who-owns-what comparison against Ch.8 (splay intro), Ch.22 (HLD),
%  Ch.16 (LCA), Ch.5 (Euler tour) and Ch.24 (centroid), the recognition
%  guide, a constraint->technique table, and notation.
%  Exports fig:c23-tree, fig:c23-taxonomy.
%  No \term{splay tree} here -- it is indexed by Ch.8 (\chref{c06}); \term the
%  NEW objects "link-cut tree", "preferred path", "represented tree" only.
% =====================================================================

\begin{chaptertheorybox}{Chapter at a Glance}
Chapter~8 showed you a \emph{splay tree}: a binary search tree that rotates the
node you just touched up to the root, so frequently-used keys float to the top and
a long run of accesses costs \(O(\log n)\) each \emph{amortised} (\secref{c06-splay}).
It closed with an explicit IOU -- "full analysis, top-down splaying, and link-cut
dynamic trees are Chapter~25."  This chapter pays all three.

A link-cut tree answers a question no static structure can: keep a forest whose
edges you \emph{add and delete online}, and after every change return an aggregate
(sum, min, max) over the unique path between two nodes, or decide whether they are
even connected.  The trick is to store the real tree (the \term{represented tree})
as a handful of splay trees (the \term{auxiliary trees}) -- one per \term{preferred
path} -- and let a single operation, \code{access}, re-shuffle which edges are
preferred so that the root-to-node path you care about is always one contiguous
splay tree you can read in \(O(\log n)\).

\textbf{You will learn to}
\begin{itemize}[nosep]
  \item run the \emph{splay engine} -- splaying by access, the zig / zig-zig /
        zig-zag steps, the potential function that makes one operation \(O(n)\)
        worst-case but \(O(\log n)\) amortised, top-down splaying with no
        recursion, and size-maintained splay trees as an order-statistics set
        (\(\secref{c23-splay}\), solved SPOJ ORDERSET); \hfill\lv{3}--\lv{4}
  \item cut a rooted forest into \emph{preferred paths}, store each as one aux
        splay keyed by depth, walk the \code{access} spine, and read the LCA of
        two nodes with two accesses and no rerooting
        (\(\secref{c23-prefer}\), solved SPOJ DYNALCA); \hfill\lv{4}--\lv{5}
  \item \emph{evert} a tree with \code{makeroot} + a reverse lazy tag, then
        \code{link}, \code{cut}, and fold a vertex monoid up the aux trees to
        answer a path aggregate in \(O(\log n)\)
        (\(\secref{c23-makeroot}\), the worked dynamic path-sum); \hfill\lv{5}
  \item push \emph{beyond paths} -- maintain a \code{vsum} accumulator for virtual
        (non-preferred) children so subtree aggregates survive an access, and see
        why connectivity is often an Euler-tour-tree or rollback-DSU job, not an
        LCT job (\(\secref{c23-beyond}\)); \hfill\lv{5}
  \item read a specification and decide whether the topology is even dynamic --
        if it is fixed, hand the job back to HLD / Euler+Fenwick / centroid
        (\(\secref{c23-boundary}\), the pick table). \hfill\lv{5}
\end{itemize}

\textbf{Prerequisites.} Chapter~8 (\chref{c06}) -- the splay tree and its
amortised behaviour, treated here as the engine we finally drive all the way.
Chapter~16 (\chref{c11}) -- \(\lca\) and the static \(\mathit{tin}\)/depth world a
link-cut tree replaces.  Chapter~22 (\chref{c16}) -- heavy-light decomposition:
the \emph{static} cousin of a preferred-path split, cited (never re-taught) in
\(\secref{c23-prefer}\).  Chapter~5 (\chref{c24}) -- the Euler tour that becomes
an Euler-tour tree in \(\secref{c23-beyond}\).  Chapter~24 (\chref{c18}) -- the
offline divide-and-conquer engine that beats an LCT for a known edge timeline.
\end{chaptertheorybox}

% ---------------- Reading paths + ladder position ----------------
\begin{chaptertheorybox}{How to Read This Chapter (three paths)}
Level badges \(\lv3\)--\(\lv6\): \(\lv3\) is splaying a node to the root of a BST
by hand and counting rotations; \(\lv{4}\) is a size-maintained splay set that
answers k-th / count-less-than, or a correct \code{access} spine on the running
forest; \(\lv{5}\) is \code{makeroot} with a reverse tag driving link/cut/path-sum,
or a subtree \code{vsum} fix-up; \(\lv{6}\) is the Summit -- a block-cut tree with
an LCT path-min riding on top of it.

\textbf{Engine path (\(\lv3\)--\(\lv{5}\), "make splaying second nature").}
\secref{c23-splay} in full -- the three splay cases, the potential, top-down
splaying, size aggregates, solved SPOJ ORDERSET -- then \secref{c23-prefer}, where
the same splay becomes the auxiliary tree of a preferred path (solved SPOJ
DYNALCA), capped by the Summit \(\probref{c23-summit}\).

\textbf{Full-LCT path (\(\lv{4}\)--\(\lv{6}\), "a dynamic forest I can query").}
\secref{c23-prefer} for \code{access} and rooted LCA, then \secref{c23-makeroot}
for evert / link / cut / path aggregate (the worked dynamic path-sum), then
\secref{c23-beyond} for the subtree \code{vsum} that makes a real competitive LCT
correct -- and the Summit, where an LCT sits on a block-cut tree.

\textbf{Boundary path (\(\lv{4}\)--\(\lv{5}\), "is this even a link-cut job?").}
\secref{c23-boundary}: the pick table that hands fixed-tree path ranges to HLD
(\chref{c16}), subtree ranges to Euler+Fenwick (\chref{c24}), offline edge
timelines to rollback-DSU / centroid (\chref{c18}), and pure connectivity to an
Euler-tour tree -- so you only pay the LCT constant when the topology truly moves.
\end{chaptertheorybox}

\begin{figure}[H]
\centering
\begin{tikzpicture}[diagram, x=8.5mm, y=6.5mm]
  \foreach \st/\lab in {1/lv1,2/lv2,3/lv3,4/lv4,5/lv5,6/lv6} {
    \fill[\lab!22, draw=\lab, line width=0.8pt] (\st-1,0) rectangle (\st,\st*0.52);
    \node[font=\scriptsize\bfseries, text=\lab] at (\st-0.5,\st*0.52-0.42) {L\st};}
  \node[annot, anchor=south west, font=\scriptsize]
    at (0.05,0.62) {prereq (Ch.8): a splay tree, rotate-to-root};
  \node[annot, anchor=south west, font=\scriptsize]
    at (1.05,1.14) {25.1 zig / zig-zig / zig-zag; potential; amortised O(log n)};
  \node[annot, anchor=south west, font=\scriptsize]
    at (2.05,1.66) {25.1 top-down splay, size aggregates, solved SPOJ ORDERSET};
  \node[annot, anchor=south west, font=\scriptsize]
    at (3.05,2.18) {25.2 preferred paths, access spine, solved SPOJ DYNALCA};
  \node[annot, anchor=south west, font=\scriptsize, text=lv5]
    at (4.05,2.70) {25.3 makeroot/link/cut/path-sum; 25.4 subtree vsum, ETT};
  \node[annot, anchor=south west, font=\scriptsize, text=lv6]
    at (5.05,3.22) {25.5 the pick table; Summit CF 487E (block-cut tree + LCT)};
\end{tikzpicture}
\caption*{Chapter~25 on the six-step ladder: splaying a node to the root by hand is
\(\lv3\); a size-maintained splay set and a correct access spine are \(\lv{4}\);
\code{makeroot} with a reverse tag and a subtree \code{vsum} are \(\lv{5}\); the
Summit (an LCT path-min on a block-cut tree) is step~6.}
\end{figure}

% ---------------- The running example ----------------
\subsection*{The Running Example: One Rooted Forest, Accessed Over and Over}

Every section in this chapter runs on the same rooted forest \(R\) on \(n=8\)
nodes, fixed once here and authoritative (if a later trace disagrees, this figure
wins).  It is a single tree rooted at \(1\); parent-to-child edges, 1-based:
\[
  1\!\to\!2,\quad 1\!\to\!3,\quad 2\!\to\!4,\quad 2\!\to\!5,\quad
  5\!\to\!6,\quad 5\!\to\!7,\quad 3\!\to\!8 .
\]
Depths from the root are \(\depth(1)=0\), \(\depth(2)=\depth(3)=1\),
\(\depth(4)=\depth(5)=\depth(8)=2\), \(\depth(6)=\depth(7)=3\).  In \(\secref{c23-prefer}\)
we mark the preferred edges (drawn solid) as the ones on the access path to node
\(7\): \(1\!\to\!2\!\to\!5\!\to\!7\) is the preferred path containing the root, and
\(\{4\},\{6\},\{3,8\}\) hang off it as separate preferred paths -- three auxiliary
splay trees plus the root's spine.

\begin{figure}[H]
\centering
\begin{tikzpicture}[diagram, scale=0.95, transform shape,
   tnode/.style={circle, draw=primaryblue!70, fill=primaryblue!8, line width=0.7pt,
                 minimum size=6.5mm, font=\scriptsize\bfseries, inner sep=1pt},
   pref/.style={circle, draw=difforange, fill=difforange!14, line width=1pt,
               minimum size=6.8mm, font=\footnotesize\bfseries, inner sep=1pt}]
  \node[pref] (n1) at (0,0) {1};
  \node[pref] (n2) at (-2.4,-1.4) {2};
  \node[tnode] (n3) at (2.4,-1.4) {3};
  \node[tnode] (n4) at (-3.8,-2.9) {4};
  \node[pref] (n5) at (-1.2,-2.9) {5};
  \node[tnode] (n8) at (2.4,-2.9) {8};
  \node[tnode] (n6) at (-2.2,-4.4) {6};
  \node[pref] (n7) at (-0.2,-4.4) {7};
  \draw[difforange, line width=1.6pt] (n1)--(n2) (n2)--(n5) (n5)--(n7);
  \draw (n1)--(n3) (n2)--(n4) (n3)--(n8) (n5)--(n6);
  \node[annot, font=\scriptsize, text=difforange] at (2.2,0.2) {preferred path to 7 (solid, orange)};
\end{tikzpicture}
\caption{The running forest \(R\) (\(n=8\), rooted at \(1\)).  The orange solid
edges \(1\!-\!2\!-\!5\!-\!7\) are the preferred path after an \code{access(7)}; the
orange nodes are one auxiliary splay tree keyed by depth.  The white subtrees
\(\{4\}\), \(\{6\}\), \(\{3,8\}\) are each their own preferred path (their own aux
tree) whose top node's \emph{path-parent} pointer aims up into the orange spine.}
\label{fig:c23-tree}
\end{figure}

% ---------------- Which job, cost chart ----------------
\subsection*{Which Job? (the question-shape chart)}

\begin{center}
\small
\renewcommand{\arraystretch}{1.25}
\setlength{\tabcolsep}{3pt}
\begin{tabular}{@{}p{0.30\linewidth} p{0.17\linewidth} p{0.29\linewidth} >{\raggedright\arraybackslash}p{0.14\linewidth}@{}}
\toprule
\textbf{Question} & \textbf{Technique} & \textbf{Why not a static structure} & \textbf{Where} \\
\midrule
order statistics on a dynamic set (k-th, count \(<x\)) & size-maintained splay & a sorted array is \(O(n)\) per insert; a BIT needs small keys & \secref{c23-splay} \\
LCA while edges link and cut online & rooted LCT \code{access} & binary lifting is rebuilt in \(O(n\log n)\) per edit & \secref{c23-prefer} \\
aggregate over the path \(u\!-\!v\) as the tree reroots & \code{makeroot}+path fold & HLD assumes a fixed root and fixed chains & \secref{c23-makeroot} \\
\code{link(u,v)}, \code{cut(u,v)}, ``are \(u,v\) connected?'' & LCT / Euler-tour tree & no static index survives a topology edit & \secref{c23-makeroot} \\
subtree aggregate on a \emph{changing} tree & LCT + virtual \code{vsum} & a Euler tour is destroyed by a cut & \secref{c23-beyond} \\
path range query, \emph{fixed} tree & \emph{go to HLD} & the LCT constant is larger; chains are simpler & \chref{c16} \\
offline edge timeline, all queries known & \emph{rollback DSU / centroid} & you never need an online structure & \chref{c18} \\
\bottomrule
\end{tabular}
\end{center}

\begin{insightbox}[One promise decides everything: make the path you need contiguous]
A splay tree is just a BST; a link-cut tree is a \emph{set} of splay trees, each
holding one preferred path of the real forest, ordered by depth.  The only real
operation is \code{access(v)}: splay \(v\) to the top of its path and repeatedly
walk up through path-parent pointers, at each step declaring the edge you came
from \emph{preferred} and the edge you leave behind \emph{unpreferred}.  After one
\code{access}, the whole root-to-\(v\) path is a single aux splay whose in-order
sequence is exactly that path -- so a path aggregate is the splay's subtree fold,
read in \(O(\log n)\).  Because every step is a splay, the amortised bound
inherits the \(O(\log n)\) you proved in Chapter~8, and the tree can be rerooted,
relinked and recut freely -- the represented tree is a \emph{view}, the aux splay
trees are the \emph{data}.
\end{insightbox}

% ---------------- Taxonomy chart ----------------
\begin{figure}[H]
\centering
\begin{tikzpicture}[diagram, scale=0.95, transform shape,
    mnode/.style={circle, draw=#1, line width=0.9pt, fill=#1!8, minimum size=9mm,
                  font=\scriptsize\bfseries, text=#1!60!black}]
  \node[mnode=primaryblue] (a) at (0,0)   {splay};
  \node[annot, font=\scriptsize, anchor=west, text width=9.8cm, align=left] at (1.7,0) {rotate the touched node to the root; zig / zig-zig / zig-zag; the potential makes a run O(log n) amortised; keep subtree sizes for order statistics (25.1, SPOJ ORDERSET)};
  \node[mnode=successgreen] (b) at (0,-1.1) {access};
  \node[annot, font=\scriptsize, anchor=west, text width=9.8cm, align=left] at (1.7,-1.1) {split the rooted forest into preferred paths, one aux splay each; access(v) makes root-to-v one contiguous splay; two accesses read the LCA, no reroot (25.2, SPOJ DYNALCA)};
  \node[mnode=difforange]   (c) at (0,-2.2) {evert};
  \node[annot, font=\scriptsize, anchor=west, text width=9.8cm, align=left] at (1.7,-2.2) {makeroot(v) = access(v) + flip a reverse lazy tag on v's splay; now v is the tree root, so link(u,v), cut(u,v) and a path fold u--v all become access-and-read (25.3, worked path-sum)};
  \node[mnode=warningred]   (d) at (0,-3.3) {virtual};
  \node[annot, font=\scriptsize, anchor=west, text width=9.8cm, align=left] at (1.7,-3.3) {a splay child is only the PREFERRED child; a subtree aggregate needs a separate vsum over the non-preferred (virtual) children, fixed up every time access detaches a solid edge (25.4); connectivity -> Euler-tour tree / DSU};
  \node[mnode=softpurple]   (e) at (0,-4.4) {pick};
  \node[annot, font=\scriptsize, anchor=west, text width=9.9cm, align=left] at (1.7,-4.4) {is the topology truly dynamic? yes -> LCT; fixed path ranges -> HLD; fixed subtree ranges -> Euler+Fenwick; offline timeline -> rollback DSU/centroid; only connectivity -> ETT (25.5)};
\end{tikzpicture}
\caption{The chapter's five-state taxonomy.  \emph{splay} is the engine; \emph{access}
turns the engine into a rooted dynamic-LCA structure; \emph{evert} adds rerooting
and so full link/cut/path-aggregate power; \emph{virtual} is the extra accumulator a
real LCT needs for subtree queries and the connectivity-vs-path split; \emph{pick}
is the boundary choosing between an LCT and its cheaper neighbours.  Each solved
problem sits on one row: ORDERSET on \emph{splay}, DYNALCA on \emph{access}, the
worked path-sum on \emph{evert}, the \code{vsum} example on \emph{virtual}, and the
Summit CF~487E on \emph{evert}+a block-cut condensation.}
\label{fig:c23-taxonomy}
\end{figure}

% ---------------- Who owns what ----------------
\subsection*{Who Owns What (so nothing is learned twice)}

\begin{center}
\small
\renewcommand{\arraystretch}{1.25}
\begin{tabular}{@{}p{0.30\linewidth} p{0.24\linewidth} p{0.40\linewidth}@{}}
\toprule
\textbf{Idea} & \textbf{Owner} & \textbf{This chapter's share} \\
\midrule
splay tree definition, rotate-to-root & Ch.~8 (\chref{c06}) & the IOU is paid: full potential analysis, top-down splaying and size aggregates (25.1); the single rotation is never re-drawn \\
static heavy-light path split & Ch.~22 (\chref{c16}) & the preferred path is the SAME heavy-path idea made dynamic; cited once in 25.2, boundary drawn in 25.5 \\
\(\lca\) by binary lifting & Ch.~16 (\chref{c11}) & 25.2's two-access LCA is an ALTERNATIVE for a moving forest, not a re-derivation \\
Euler tour, flatten to array & Ch.~5 (\chref{c24}) & becomes an Euler-tour \emph{tree} for connectivity in 25.4; tin/tout never re-taught \\
offline divide-and-conquer on a timeline & Ch.~24 (\chref{c18}) & 25.5 routes a known edge timeline to rollback-DSU / centroid instead of a live LCT \\
dynamic connectivity (only) & Ch.~17 DSU (\chref{c28}) & if you never need a path value, a (rollback) DSU beats an LCT; forward-cited in 25.4 \\
compressed tree of a marked node set & Ch.~26 (\chref{c30}) & explicitly out; virtual trees are the next chapter \\
\bottomrule
\end{tabular}
\end{center}

% ---------------- Recognition guide ----------------
\begin{recognitionbox}[Pattern card -- "the tree itself changes under my queries"]
\reckey{statements with an online \emph{edge edit} on a forest: ``\code{link(u,v)}
/ \code{cut(u,v)} / query the path sum from \(u\) to \(v\), all interleaved'',
``support rerooting the tree, then a path aggregate'', ``after adding or deleting a
road, is \(u\) still reachable from \(v\)? and what is the cheapest city on the
route?'', ``maintain a dynamic forest and answer a value that lives on the unique
path between two live nodes''.  The tell is a \emph{topology edit} (an edge appears
or vanishes) interleaved with a \emph{path} query -- if the edges are fixed, you
wanted HLD; if you only need connectivity, you wanted DSU.}
\recuse{
  \item edges link and cut online \emph{and} you query a value along the current path -- full LCT with \code{makeroot} (\(\secref{c23-makeroot}\));
  \item a rooted forest where only links/cuts change and you need LCA -- rooted LCT, no evert (\(\secref{c23-prefer}\));
  \item a dynamic ordered set / sequence needing k-th or rank -- size-maintained splay, the engine alone (\(\secref{c23-splay}\)).
}
\recnotuse{
  \item the tree is \emph{fixed} and you query/update a path range -- HLD (\chref{c16}), smaller constant;
  \item the tree is \emph{fixed} and you query a subtree -- Euler tour + Fenwick (\chref{c24});
  \item all edge edits are \emph{known in advance} (offline) -- rollback DSU or segment-tree-on-time, no online structure;
  \item you only need \emph{connectivity}, never a path value -- (Euler-tour) DSU is simpler (\(\secref{c23-beyond}\));
  \item a \emph{subset of marked nodes} compressed to its Steiner tree -- virtual trees (\chref{c30}).
}
One-line test: \emph{an edge appears or disappears between two queries AND the
answer depends on the current path -- then a link-cut tree; otherwise a static
engine is faster.}
\end{recognitionbox}

% ---------------- Constraint -> technique ----------------
\subsection*{Reading the Constraints (which technique the limits want)}

\begin{center}
\small
\renewcommand{\arraystretch}{1.25}
\begin{tabular}{@{}p{0.34\linewidth} p{0.40\linewidth} p{0.20\linewidth}@{}}
\toprule
\textbf{Constraint signature} & \textbf{What it rules in/out} & \textbf{Technique} \\
\midrule
dynamic set, keys \(|x|\le10^9\), k-th / rank queries & a BIT cannot index huge keys; a splay holds them and carries sizes & \secref{c23-splay} \\
link/cut a forest, ask LCA, \(N,M\le10^5\) & binary lifting rebuilds per edit; \code{access} answers in \(O(\log n)\) & \secref{c23-prefer} \\
\code{link}, \code{cut}, path sum/min interleaved online & only a rerootable aux-splay forest survives; LCT & \secref{c23-makeroot} \\
subtree sum on a tree that also link/cuts & preferred children miss the off-path subtrees; add a virtual \code{vsum} & \secref{c23-beyond} \\
edge insert/delete timeline given \emph{offline}, connectivity only & no online structure needed; rollback DSU / D\&C (\chref{c18}) & \secref{c23-boundary} \\
\bottomrule
\end{tabular}
\end{center}

% ---------------- Notation ----------------
\subsection*{Notation Used Only Here}

The represented forest is \(R=(V,E)\), \(|V|=n\), each tree rooted (until an
\code{evert} changes which node is root).  A node's \emph{aux} children are the
at-most-two children in its splay tree (the \emph{preferred} path neighbours, one
shallower and one deeper); its \emph{path-parent} is a single pointer from the top
of a preferred path to the node just above it in the represented tree.  A node
\(x\) is an \term{aux root} (\code{is\_root(x)}) when it is not a child in any
splay tree -- equivalently its parent, if any, treats it as a path-parent, not an
aux child.  \code{access(v)} detaches \(v\)'s right (deeper) preferred child and
repeatedly climbs path-parents, so after it \(v\)'s splay holds the whole current
root-to-\(v\) path in depth order.  \code{makeroot(v)} = \code{access(v)} then a
reverse lazy tag on \(v\)'s splay, making \(v\) the represented root.  The
aggregate of an aux splay node \(x\) is \(\mathit{sum}(x)\) (or \(\min\)) over its
splay subtree -- i.e.\ over the preferred path segment \(x\) represents; a
\emph{virtual child} is a represented child not currently on \(x\)'s preferred
path, and their values are held in \(\mathit{vsum}(x)\) (\(\secref{c23-beyond}\)).
\(\depth(\cdot)\) is depth in the current represented tree.  Every section runs on
\(R\) above; if a table here and a trace there disagree, this opener is
authoritative.

% =====================================================================
%  25.1  The Splay Engine -- rotations, the potential, top-down splaying,
%  and size-maintained splay trees as an order-statistics set.
%  Exports: sec:c23-splay  def:c23-splay  thm:c23-splay-amort
%           fig:c23-splaycases  lst:c23-splay-rot  lst:c23-topdown
%           prob:c23-orderset  lst:c23-orderset
%  Solved SPOJ ORDERSET (Order statistic set), live-verified via the SPOJ
%  problem page (tags #tree #binary-search; dynamic set, Q <= 200000,
%  |x| <= 1e9). Official sample read live:
%    8 ops: I -1, I -1, I 2, C 0, K 2, D -1, K 1, K 2  ->  1, 2, 2, invalid.
%  Hand-replayed here: S={-1,2}; C 0 -> 1; K 2 -> 2; D -1 -> S={2}; K 1 -> 2;
%  K 2 -> invalid (only one element). Matches the official output exactly.
%  DISCHARGES the Ch.8 sec:c06-splay forward-cite ("full analysis, top-down
%  splaying ... Chapter 25"); the single rotation is NOT re-drawn.
% =====================================================================
\section{The Splay Engine: Balance by Locality}\label{sec:c23-splay}

\begin{previouslybox}[Ch.~8 built the splay tree and left the analysis running]
Chapter~8 (\(\secref{c06-splay}\)) introduced the \term{splay tree}: a binary
search tree that, every time you touch a node, rotates it all the way to the root
(\(\figref{c06-splayfig}\)). It gave the three rotation names -- zig, zig-zig,
zig-zag -- and promised that a long run of operations costs \(O(\log n)\) each
\emph{on average}, then explicitly deferred \emph{"the full analysis, top-down
splaying, and link-cut dynamic trees"} to this chapter. That IOU is cashed here:
\(\thmref{c23-splay-amort}\) proves the amortised bound with a potential function,
the top-down variant replaces the recursive splay with a single downward pass, and
subtree sizes turn the engine into an order-statistics set -- the exact tool
SPOJ~ORDERSET wants. Only then, in \(\secref{c23-prefer}\), do we hand this engine
to the link-cut tree.
\end{previouslybox}

\begin{topicintro}
A splay tree is a plain BST plus one habit: \code{access(x)} does the ordinary
search for \(x\), then \emph{splays} \(x\) to the root -- repeatedly rotating \(x\)
up two levels at a time using zig / zig-zig / zig-zag until it sits at the top.
Three things fall out of that habit. (1) Recently- and frequently-accessed keys
cluster near the root, so hot data is fast. (2) Even though one access can cost
\(\Theta(n)\) (a chain of cold nodes), the \emph{amortised} cost per operation is
\(O(\log n)\) -- proved with the potential \(\Phi=\sum_v \log \size(v)\)
(\(\thmref{c23-splay-amort}\)). (3) Because splaying only needs rotations, any
aggregate you can push up a tree (subtree \(\size\), sum, min) rides along for
free, giving an implicit, order-statistics set with no rotations-of-aggregates
surprise.

\textbf{Prerequisites.} Chapter~8 (\(\secref{c06-splay}\)) -- the BST shape, the
single rotation, and the names zig / zig-zig / zig-zag. This section reuses that
vocabulary and adds only the analysis and the two new payload features (top-down
splaying, subtree sizes); the mechanics of one rotation are never re-drawn.

\textbf{After this section you can}
\begin{itemize}[nosep]
  \item splay a node to the root by hand, naming each zig / zig-zig / zig-zag; \hfill\lv{3}
  \item prove the amortised \(O(\log n)\) bound from \(\Phi=\sum\log\size\); \hfill\lv{4}
  \item write top-down splaying -- one downward pass, no parent pointers, no recursion; \hfill\lv{4}
  \item maintain subtree sizes in a splay and solve SPOJ~ORDERSET (k-th, count-\(<x\))
        on keys up to \(10^9\). \hfill\lv{4}
\end{itemize}
\end{topicintro}

\begin{recognitionbox}[Pattern card -- "a dynamic ordered set with rank/select"]
\reckey{statements that keep a \emph{set of numbers} under \code{INSERT}/\code{DELETE}
and ask \code{K-TH} (the \(k\)-th smallest) or \code{COUNT} (how many are \(<x\)),
with \emph{huge keys} (\(|x|\le 10^9\)) so an array-indexed BIT is impossible and
coordinate-compression is awkward because inserts and deletes arrive online. The
tell is order-statistics on a live set -- rank and select -- where the operations
themselves, not the key range, set the budget.}
\recuse{
  \item dynamic set, k-th smallest / count-less-than interleaved with insert-delete -- size-maintained splay (\(\secref{c23-splay}\));
  \item as the auxiliary tree of a preferred path -- the same engine, no sizes, in the LCT (\(\secref{c23-prefer}\));
  \item implicit sequence (split/merge by position) -- splay with subtree sizes as an index (\(\secref{c23-splay}\)).
}
\recnotuse{
  \item keys are small and fixed, only prefix counts -- a Fenwick tree (\chref{c14}) is simpler;
  \item the set never shrinks and you only need membership -- \code{std::set} / hashing;
  \item you need the \(k\)-th but the set is \emph{static} -- sort once, binary search.
}
One-line test: \emph{rank or select on a set that grows and shrinks online, with
keys you cannot index into an array -- then a size-maintained splay tree.}
\end{recognitionbox}

% ---------------------------------------------------------------
\subsection{Splaying, and why its average is logarithmic}
% ---------------------------------------------------------------

\begin{definition}{Splaying a node to the root}{c23-splay}
Let \(x\) be a node with parent \(p\) and grandparent \(g\) (either may be absent).
\term{Splaying} \(x\) repeats, until \(x\) is the root:
\begin{itemize}[nosep]
  \item \textbf{zig} (\(p\) is the root): one rotation at \(x\)'s parent; \(x\) becomes root.
  \item \textbf{zig-zig} (\(x,p\) on the \emph{same} side of \(g\)): rotate at \(g\),
        then at \(p\) -- both rotations same direction, \(x\) climbs two levels and
        \(p\) stays below \(x\).
  \item \textbf{zig-zag} (\(x,p\) on \emph{opposite} sides of \(g\)): rotate at \(p\),
        then again at \(p\)'s new child -- \(x\) climbs two levels and takes \(p\)'s old spot.
\end{itemize}
The BST order is preserved throughout; only shape changes.
\end{definition}

\begin{figure}[H]
\centering
\begin{tikzpicture}[diagram, scale=0.72, transform shape,
   snode/.style={circle, draw=primaryblue!70, fill=primaryblue!8, line width=0.7pt,
                 minimum size=6mm, font=\scriptsize\bfseries, inner sep=1pt},
   hl/.style={circle, draw=difforange, fill=difforange!16, line width=1pt,
              minimum size=6.4mm, font=\scriptsize\bfseries, inner sep=1pt}]
  % zig: g absent, p root, x child  -> x root
  \node[snode] (zp) at (0,0) {p}; \node[hl] (zx) at (1.1,-1.2) {x};
  \draw (zp)--(zx);
  \node[annot, font=\scriptsize, anchor=south] at (0.5,0.5) {zig};
  % zig-zig: g-p-x same side
  \node[snode] (gz) at (-6.4,0) {g}; \node[snode] (pz) at (-5.3,-1.2) {p}; \node[hl] (xz) at (-4.2,-2.4) {x};
  \draw (gz)--(pz)--(xz);
  \node[annot, font=\scriptsize, anchor=south] at (-5.3,0.5) {zig-zig (x,p same side: rot g, rot p)};
  % zig-zag: g, p, x opposite sides
  \node[snode] (ga) at (5.0,0) {g}; \node[snode] (pa) at (3.9,-1.2) {p}; \node[hl] (xa) at (5.0,-2.4) {x};
  \draw (ga)--(pa) (pa)--(xa);
  \node[annot, font=\scriptsize, anchor=south] at (4.5,0.5) {zig-zag (rot p twice)};
\end{tikzpicture}
\caption{The three splay cases around a grandparent \(g\), parent \(p\), target
\(x\) (\(x\) in orange). \emph{zig} ends the splay (no grandparent). \emph{zig-zig}
and \emph{zig-zag} each lift \(x\) two levels; the choice -- same side versus
opposite side -- is exactly what keeps the tree self-correcting and makes the
amortised bound work (\(\thmref{c23-splay-amort}\)). Names and mechanics are
Ch.~8's (\(\figref{c06-splayfig}\)); here they are only relabelled for reference.}
\label{fig:c23-splaycases}
\end{figure}

\begin{theorem}{Splaying is \(O(\log n)\) amortised per operation}{c23-splay-amort}
Over any sequence of \(m\) insertions, deletions and searches in a splay tree on
\(n\) keys, the total time is \(O(m\log n)\); the amortised cost of one operation
is \(O(\log n)\). A single operation may still cost \(\Theta(n)\).
\end{theorem}
\begin{proofbox}
Define the \emph{rank} \(\rho(v)=\log \size(v)\) (size = nodes in \(v\)'s subtree)
and the potential \(\Phi=\sum_{v}\rho(v)\). Charge a step's \emph{amortised} cost as
\(\text{actual} + \Delta\Phi\). A single rotation at a node \(t\) (reparenting a
subtree of size \(S\) between two configurations of sizes \(A,B\)) changes only the
ranks of \(t\) and its child; bounding that change by \(\Delta\Phi_{\text{rot}}\le
2(\rho_{\text{new}}(t)-\rho_{\text{old}}(t))\) and summing the telescoping terms
along the splay path \(x=v_0,v_1,\dots,v_k=\text{root}\) gives the access of \(x\) a
total amortised cost at most
\[
  3\sum_{i}\bigl(\rho(v_{i,\text{new}})-\rho(v_{i,\text{old}})\bigr) + O(1)
  \;\le\; 3\bigl(\rho(\text{root})-\rho(x)\bigr) + O(\log n) \;=\; O(\log n),
\]
because \(\rho(\text{root})=\log n\) after the splay and the series telescopes. The
zig-zig / zig-zag pairing is what forces the factor \(3\) (a plain single-step
rotation gives only \(O(\log^2 n)\)); this is the standard Sleator--Tarjan
potential argument. Insert and delete are a search-and-splay plus one splice, so
they inherit the same bound; summing \(m\) operations telescopes \(\Phi\) to
\(O(m\log n)\) total. \(\Box\)
\end{proofbox}

\begin{complexitybox}[The engine]
\textbf{Per operation} \(\bigO{\log n}\) amortised (\(\thmref{c23-splay-amort}\)),
\(\Theta(n)\) worst case. \textbf{Space} \(\bigO{n}\). Maintaining a subtree
aggregate such as \(\size\) adds \(O(1)\) work per rotation (recompute the two
affected nodes), so it costs nothing asymptotically -- this is why sizes, sums and
lazy tags ride a splay for free.
\end{complexitybox}

\begin{lstlisting}[caption={Bottom-up splay: rotate x up two levels at a time. grand(x) is x's grandparent; the zig-zig vs zig-zag split is the whole art. Assumes parent pointers. ASCII only}, label={lst:c23-splay-rot}]
// ch(x,0)=left, ch(x,1)=right ; p(x)=parent ; rot(x) rotates x with its parent
void rot(Node* x){
    Node* p = par[x]; Node* g = par[p];
    int d = (ch[p][1] == x);            // x is p's right child?
    ch[p][d] = ch[x][!d]; if (ch[x][!d]) par[ch[x][!d]] = p;
    ch[x][!d] = p; par[p] = x; par[x] = g;
    if (g) ch[g][ (ch[g][1] == p) ? 1 : 0 ] = x;
    pull(p); pull(x);                    // refresh sizes/sums bottom-up
}
void splay(Node* x){
    while (par[x]) {
        Node* p = par[x]; Node* g = par[p];
        if (!g) rot(x);                          // zig
        else if ((ch[g][0]==p) == (ch[p][0]==x)) { rot(p); rot(x); }  // zig-zig
        else { rot(x); rot(x); }                 // zig-zag
    }
    root = x;
}
\end{lstlisting}

\begin{lstlisting}[caption={Top-down splaying: one downward pass, no recursion, no parent pointers. Keys smaller than the walked-down spine accumulate on tree L, larger on tree R; x is re-attached at the end as the new root. ASCII only}, label={lst:c23-topdown}]
// Splay key `key` to the root in one pass. l/r are dummy roots of the two chains.
Node* topdown_splay(Node* t, int key){
    Node L, R; L.right = &L; R.left = &R;        // local sentinels
    Node *l = &L, *r = &R;
    while (t) {
        if (key < t->key) {                       // go left: t and its right subtree >= key
            Node* tmp = t->right; t->right = r->left; r->left = t; r = t; t = tmp;
        } else if (key > t->key) {                // go right: t and its left subtree <= key
            Node* tmp = t->left; t->left = l->right; l->right = t; l = t; t = tmp;
        } else break;                             // found t; assemble below
    }
    if (!t) { root = &L, root = &R; return; }     // key absent
    // assemble: t's left  gets L's right-chain, t's right gets R's left-chain
    l->right = t->left; t->left = L.right;
    r->left  = t->right; t->right = R.left;
    pull_down(t);                                  // recompute sizes along the two chains
    root = t;
}
\end{lstlisting}

\begin{insightbox}[Top-down beats bottom-up where recursion is the enemy]
Bottom-up splaying (\(\coderef{c23-splay-rot}\)) recurses up through parent
pointers and then back down to fix sizes -- two passes and \(O(\log n)\) stack
(amortised) that can degenerate to \(O(n)\) stack on a cold chain, risking a
recursion-overflow on a strict judge. Top-down splaying (\(\coderef{c23-topdown}\))
walks \emph{down} once, peeling the searched path into two side chains \(L\) and
\(R\), and re-attaches them under \(x\). Same \(O(\log n)\) amortised bound, no
parent array, no recursion -- the variant you actually ship. Ch.~8 flagged it as
deferred; here it is written out.
\end{insightbox}

% ---------------------------------------------------------------
\subsection{Problem Lab: SPOJ ORDERSET -- rank and select on a live set}
% ---------------------------------------------------------------

\solvedproblem{c23-orderset}{Order Statistic Set}{SPOJ ORDERSET}
  {https://www.spoj.com/problems/ORDERSET/}{SPOJ classic L4}
  {size-maintained splay; k-th = select by size; count-<x = BST walk with sizes}

\begin{problemstatement}[Order statistic set (statement and sample verified against the live SPOJ page)]
Maintain a set \(S\) of integers (duplicates ignored) under four operations.
\(\texttt{I }x\): insert \(x\) if absent. \(\texttt{D }x\): delete \(x\) if present.
\(\texttt{K}\): return the \(k\)-th smallest element of \(S\), or the word
\texttt{invalid} if \(|S|<k\). \(\texttt{C }x\): return the number of elements
strictly less than \(x\). \(Q\le 2\cdot 10^5\) operations, \(|x|\le 10^9\).

\textbf{Example (official).}
\begin{verbatim}
Input:              Output:
8                   1
I -1                2
I -1                2
I 2                 invalid
C 0
K 2
D -1
K 1
K 2
\end{verbatim}
\end{problemstatement}

\begin{attemptgate}[Attempt it yourself first -- \lv{4}, give yourself 30 minutes]
Ignore the \(10^9\) keys and the delete for one moment: how would a BST answer
\code{K}? \peek{to print the \(k\)-th smallest you must know, at each node, how many
keys sit in its left subtree.}
\nudge{if every node stores \(\size(\text{left})+1\) = its own 1-based rank among its
subtree, selecting the \(k\)-th is a single descent: go left if \(k\le\) left-size,
else skip them and recurse right for rank \(k-\text{leftsize}-1\).}
\fullhint{Store a subtree size in each splay node. \code{K}: descend by size as above
(this is why you splay -- you can pick any element as root since the descent only
uses sizes). \code{C x}: BST-search \(x\), accumulating \(\size(\text{left})+1\) each
time you branch right; splay the last node seen and return the count.
\(\texttt{I}\): standard splay insert (splay the inserted node up). \(\texttt{D }x\):
splay \(x\) to the root, then splice its left and right subtrees together (max of
left as new root, or simply detach) -- set semantics, so duplicates never appear.
Keys up to \(10^9\) are stored as-is; no coordinate compression needed.}
\end{attemptgate}

\solstep{Understanding}
ORDERSET is the splay engine with \emph{one} aggregate -- a subtree size -- doing all
four jobs. There is no tree topology to track (it is a plain ordered set), so the
entire problem is "keep \(\size\) correct through every rotation and use it to select
and rank." That makes it the honest warm-up before the LCT: you meet top-down
splaying, size bookkeeping, and duplicate-aware insert/delete in a setting where a
bug shows up as a wrong number, not a silent pointer corruption.

\solstep{Key Observations}
\begin{observation}
With \(\size(v)=1+\size(\text{left})+\size(\text{right})\) maintained through
rotations, the \code{K} (select) and \code{C} (rank) queries are pure descents:
select compares \(k\) against \(\size(\text{left})+1\); rank adds
\(\size(\text{left})+1\) whenever it turns right. Both are \(O(\log n)\) amortised
because each ends in a splay.
\end{observation}
\begin{observation}
Set semantics kill duplicates: on \(\texttt{I }x\) with \(x\) already present, do
nothing. The cleanest way to test presence is to splay toward \(x\); if the node
you reach already holds \(x\), it is a duplicate. This keeps \(|S|\) honest, which
\(\thmref{c23-splay-amort}\)'s size-based potential counts on.
\end{observation}

\solstep{The Approach Ladder}
\begin{center}\renewcommand{\arraystretch}{1.25}\setlength{\tabcolsep}{3pt}
\begin{tabular}{@{}c p{0.26\linewidth} p{0.30\linewidth} l c@{}}
\toprule
\textbf{\#} & \textbf{Approach} & \textbf{Idea} & \textbf{Time} & \textbf{Verdict}\\
\midrule
1 & \lv{2}\ sorted vector & keep array sorted, binary search each op & \(O(n)\) insert & \tle \\
2 & \lv{3}\ Fenwick on compressed keys & coordinate-compress all \(x\) offline & \(O(\log n)\) & \textcolor{warningred}{\ding{55}\ \textbf{MLE/awkward}} \\
3 & \lv{4}\ \textbf{size-maintained splay} & select/rank by \(\size\) (\(\coderef{c23-topdown}\)) & \(\bigO{\log n}\) amort & \ac \\
\bottomrule
\end{tabular}\end{center}
Rung~2 works only because you can read all queries first; it is fragile if you must
answer online and is the wrong instinct to bring to an LCT. Rung~3 stores the huge
keys directly and answers rank/select with the size field alone.

\solstep{Solution}
Build a top-down splay keyed on the value, each node carrying
\(\size=1+\size_L+\size_R\). Map the four commands: \texttt{I} = splay-insert,
\texttt{D} = splay-then-splice, \texttt{K} = select-by-size then splay, \texttt{C} =
rank-then-splay. Print \texttt{invalid} when \(k>|S|\).

\solstep{Implementation}
\begin{lstlisting}[caption={SPOJ ORDERSET with a bottom-up size-maintained splay (swap in lst:c23-topdown if recursion depth is a worry). Node 0 is the null sentinel with size 0. ASCII only}, label={lst:c23-orderset}]
#include <bits/stdc++.h>
using namespace std;
const int MAXN = 200005;
int ch[MAXN][2], par[MAXN], sz[MAXN], key[MAXN], root=0, tot=0;
int isNew(int k){ int id=++tot; key[id]=k; sz[id]=1; ch[id][0]=ch[id][1]=par[id]=0; return id; }
void pull(int x){ if(x) sz[x]=1+sz[ch[x][0]]+sz[ch[x][1]]; }
bool isRoot(int x){ int p=par[x]; return p==0 || (ch[p][0]!=x && ch[p][1]!=x); }
void rot(int x){
    int p=par[x], g=par[p], d=(ch[p][1]==x);
    if(!isRoot(p)) ch[g][ch[g][1]==p]=x;
    par[x]=g;
    ch[p][d]=ch[x][!d]; if(ch[x][!d]) par[ch[x][!d]]=p;
    ch[x][!d]=p; par[p]=x; pull(p); pull(x);
}
void splay(int x){
    while(!isRoot(x)){ int p=par[x], g=par[p];
        if(isRoot(p)) rot(x);
        else if((ch[g][0]==p)==(ch[p][0]==x)){ rot(p); rot(x); }
        else { rot(x); rot(x); } }
    root=x;
}
int findNode(int k){ int u=root;
    while(u && key[u]!=k) u = (k<key[u])?ch[u][0]:ch[u][1]; return u; }
void insert(int k){
    if(!root){ root=isNew(k); return; }
    int u=root, p=0;
    while(u){ p=u; u = (k==key[u]) ? 0 : (k<key[u]?ch[u][0]:ch[u][1]); }
    if(k==key[p]){ splay(p); return; }          // duplicate: set ignores it
    int nd=isNew(k); ch[p][k>key[p]]=nd; par[nd]=p;
    for(int a=nd;a;a=par[a]) pull(a); splay(nd);
}
int kth(int k){                                // k-th smallest, 1-based
    if(sz[root]<k) return -1;                  // -> "invalid"
    int u=root;
    while(true){ int L=ch[u][0];
        if(k<=sz[L]) u=L;
        else if(k==sz[L]+1) break;
        else { k-=sz[L]+1; u=ch[u][1]; } }
    splay(u); return key[u];
}
int countLess(int k){                          // #elements strictly < k
    int u=root, ans=0;
    while(u){ if(key[u]>=k) u=ch[u][0];
        else { ans += sz[ch[u][0]]+1; u=ch[u][1]; } }
    if(u==0){ /* not found, ans already complete */ }
    return ans;
}
void erase(int k){ int u=findNode(k); if(!u) return; splay(u);
    int l=ch[u][0], r=ch[u][1]; ch[u][0]=ch[u][1]=0;
    if(l){ par[l]=0; if(r){ int m=l; while(ch[m][1]) m=ch[m][1];
        ch[m][1]=r; par[r]=m; pull(m);} root=l; }
    else root=r; if(root) par[root]=0; }
int main(){
    int Q; scanf("%d",&Q);
    while(Q--){ char c; int x; scanf(" %c",&c);
        if(c=='I'){ scanf("%d",&x); insert(x); }
        else if(c=='D'){ scanf("%d",&x); erase(x); }
        else if(c=='C'){ scanf("%d",&x); printf("%d\n", countLess(x)); }
        else if(c=='K'){ scanf("%d",&x); int v=kth(x);
            if(v==-1) printf("invalid\n"); else printf("%d\n", v); } }
}
\end{lstlisting}

\solstep{Trace}
\begin{tracebox}[Official sample: sizes drive every answer]
Ops on \(S\) (shown sorted, with the size a node reports):\\
\(\texttt{I }-1\): \(S=\{-1\}\). \(\texttt{I }-1\): duplicate, ignored, \(S=\{-1\}\).
\(\texttt{I }2\): \(S=\{-1,2\}\), \(\size(\text{root})=2\).\\
\(\texttt{C }0\): walk for \(<0\): \(-1<0\) so add \(\size(\text{left of }{-1})+1=1\),
go right; \(2\not<0\) go left, hit null \(\Rightarrow\) answer \(\boxed{1}\).\\
\(\texttt{K }2\): \(\size(\text{root})=2\ge2\); descend, skip \(\{-1\}\) taking right,
\(k=1\) at node \(2\) \(\Rightarrow\) answer \(\boxed{2}\).\\
\(\texttt{D }-1\): splay \(-1\) to root, splice its right subtree up \(\Rightarrow\)
\(S=\{2\}\).\\
\(\texttt{K }1\): \(1\)-th smallest of \(\{2\}\) \(\Rightarrow\) answer \(\boxed{2}\).\\
\(\texttt{K }2\): \(|S|=1<2\) \(\Rightarrow\) print \(\boxed{\texttt{invalid}}\).\\
Outputs \(1,2,2,\texttt{invalid}\) -- matching the official sample exactly.
\end{tracebox}

\begin{warningbox}[The size field is the whole data structure]
Every ORDERSET bug is a \(\size\) bug. (1) Forgetting \code{pull} on \emph{both}
nodes touched by a rotation (parent \emph{and} grandparent chain) makes
\code{K}/\code{C} return stale ranks -- wrong answer that looks random. (2) On
\texttt{D}, splicing without first splaying \(x\) to the root leaves the children
pointers inconsistent. (3) In \code{countLess}, adding \(\size_L+1\) when you go
\emph{left} instead of right double-counts -- the count comes out too large. Verify
each with the tiny sample: the trace above is the reference oracle.
\end{warningbox}

\begin{hardeningbox}
\begin{itemize}[nosep,leftmargin=1.4em]
  \item \textbf{Duplicates.} ORDERSET is a \emph{set}. Guard \texttt{I} by the
        search result (\(\text{key}[p]==x \Rightarrow\) splay and return); a stray
        duplicate inflates \(\size\) and breaks \code{K}.
  \item \textbf{Null sentinel.} Let node \(0\) be the shared null with
        \(\text{sz}[0]=0\); then \code{pull} needs no null checks and
        \(\text{sz}[L]+1\) is always defined.
  \item \textbf{Recursion depth.} The bottom-up \code{splay}/\code{pull} loop is
        iterative; if you recurse to fix sizes, switch to \(\coderef{c23-topdown}\)
        or an explicit stack -- a cold chain is \(\Theta(n)\) deep.
  \item \textbf{\code{invalid} is lowercase.} Print exactly \texttt{invalid} when
        \(k>|S|\), not \texttt{Invalid} or \texttt{-1}.
\end{itemize}
\end{hardeningbox}

\begin{recallbox}[25.1 in one breath]
\emph{Splay = search, then rotate the touched node to the root by zig / zig-zig /
zig-zag. The potential \(\Phi=\sum\log\size\) makes a run \(O(\log n)\) amortised
(\(\thmref{c23-splay-amort}\)); top-down splaying (\(\coderef{c23-topdown}\)) does it in
one downward pass. Add one aggregate -- subtree \(\size\) -- and the same engine
answers k-th and count-\(<x\) on a live set of \(10^9\)-keys (SPOJ~ORDERSET): select
by size, rank by accumulating left-sizes. This size-carrying splay is exactly the
auxiliary tree the link-cut tree will drive in \(\secref{c23-prefer}\).}
\selftest{
  \item Which three rotations, and what decides zig-zig vs zig-zag? \emph{(whether \(x\) and \(p\) are on the same or opposite sides of \(g\), Def.~\defref{c23-splay}.)}
  \item Give the potential that yields the \(O(\log n)\) amortised bound. \emph{(\(\Phi=\sum_v \log\size(v)\), \(\thmref{c23-splay-amort}\).)}
  \item In ORDERSET, how does \code{K} locate the \(k\)-th smallest using only sizes? \emph{(compare \(k\) with \(\size_L+1\); recurse left/right. Observation~1.)}
}
\end{recallbox}

% =====================================================================
%  25.2  Preferred Paths and Access -- the rooted link-cut tree.
%  Exports: sec:c23-prefer  def:c23-prefer  def:c23-access  thm:c23-access-lca
%           fig:c23-aux  lst:c23-access  lst:c23-lct-core
%           prob:c23-dynalca  lst:c23-dynalca
%  Solved SPOJ DYNALCA (Dynamic LCA), live-verified via the SPOJ problem
%  page (forest of N<=1e5 single-vertex trees; queries link A B / cut A /
%  lca A B; M<=1e5). Official sample read live:
%    5 9: lca 1 1, link 1 2, link 3 2, link 4 3, lca 1 4, lca 3 4, cut 4,
%    link 5 3, lca 1 5  ->  1, 2, 3, 2.
%  Hand-replayed in the tracebox against the rooted-LCT algorithm here.
%  Back-cites Ch.8 splay (engine) and Ch.22 HLD (static heavy-path cousin).
% =====================================================================
\section{Preferred Paths and the Access Operation}\label{sec:c23-prefer}

\begin{previouslybox}[Ch.~22 split a \emph{fixed} tree into heavy paths; do it on the fly]
Heavy-light decomposition (\(\secref{c16-decomp}\)) cuts a tree into
\(O(\log n)\) heavy paths so any root-to-node path meets only \(O(\log n)\) of
them -- but it freezes the split in an array and cannot survive an edge edit. The
link-cut tree keeps the same philosophy with one crucial change: the choice of
"which child is preferred" is made \emph{lazily, by the access history}, and can
flip in \(O(\log n)\) whenever you touch a different node. The result is the same
"one root-to-node path is contiguous" power, now available on a forest whose edges
link and cut. The engine under it is the Chapter~8 splay tree you just mastered;
here each preferred path becomes one auxiliary splay, keyed by depth.
\end{previouslybox}

\begin{topicintro}
A \term{represented tree} (the real forest) is stored as a collection of
\term{auxiliary trees} -- splay trees, one per \term{preferred path}. A node's
\emph{preferred child} is the single child whose edge is currently "solid"; maximal
chains of solid edges are the preferred paths, and each is one aux splay whose
in-order (left-to-right) sequence is the path from its shallow top to its deep
bottom. Two pointers describe everything: a node's \emph{aux} children
(\(\mathit{ls},\mathit{rs}\)) are the neighbours on its own preferred path, and its
\emph{path-parent} (\(\mathit{fa}\) when not an aux child) points from the top of a
preferred path to the node one step above it in the real tree. The single engine
operation is \code{access(v)}: it makes the whole root-to-\(v\) path preferred, so
that afterwards \(v\)'s aux splay holds exactly that path.

\textbf{Prerequisites.} \(\secref{c23-splay}\) -- splaying and \code{is\_root}.
Chapter~22 (\(\secref{c16-decomp}\)) -- the static heavy-path picture being made
dynamic. Chapter~16 (\chref{c11}) -- \(\lca\), which we now compute without any
lifting table. This section needs no \code{makeroot}: the forest stays rooted.

\textbf{After this section you can}
\begin{itemize}[nosep]
  \item run \code{access} on the running forest \(R\) and name the aux trees and
        path-parent pointers before/after; \hfill\lv{4}
  \item test aux-root status with \code{is\_root} and read the represented root
        with \code{findroot}; \hfill\lv{4}
  \item compute the LCA of two nodes with one \code{access} plus a path-parent
        climb and no rerooting (\(\thmref{c23-access-lca}\)); \hfill\lv{5}
  \item solve SPOJ~DYNALCA end to end (link / cut / lca on a moving forest). \hfill\lv{5}
\end{itemize}
\end{topicintro}

\begin{definition}{Preferred path, auxiliary tree, path-parent}{c23-prefer}
In the link-cut representation of a rooted forest, each node has at most one
\term{preferred child}; the maximal chain of preferred (solid) edges through a node
is its \term{preferred path}, stored as one \term{auxiliary splay tree} ordered by
depth (shallower to the left, deeper to the right). The top node of a preferred
path stores a \term{path-parent} pointer \(\mathit{fa}\) to the real-tree parent of
the path's top; this pointer is \emph{not} an aux-child link, which is exactly what
\code{is\_root(x)} detects: \(x\) is an aux root iff its \(\mathit{fa}\) (if any)
does not name \(x\) as a left/right child.
\end{definition}

\begin{definition}{The access operation}{c23-access}
\code{access(v)} makes the root-to-\(v\) path preferred. Walking from \(v\) upward
through path-parent pointers, at each node \(u\) it: splays \(u\) to the top of its
aux tree; replaces \(u\)'s right child (the previously-preferred, deeper segment,
now \emph{unpreferred}) by the piece built so far; refreshes \(u\)'s aggregate; and
moves up. After the loop the final splay leaves \(v\) as the root of a single aux
tree containing the entire root-to-\(v\) path.
\end{definition}

\begin{lstlisting}[caption={The LCT core: is-root test, push-up, rotate, splay, access, findroot. Aux children are ls/rs; fa doubles as path-parent for an aux root. NO makeroot/reverse tag in this rooted section. ASCII only}, label={lst:c23-lct-core}]
int ch[MAXN][2], fa[MAXN];
bool isRoot(int x){ int f=fa[x]; return f==0 || (ch[f][0]!=x && ch[f][1]!=x); }
void rot(int x){
    int p=fa[x], g=fa[p], d=(ch[p][1]==x);
    if(!isRoot(p)) ch[g][ch[g][1]==p]=x;
    fa[x]=g;
    ch[p][d]=ch[x][!d]; if(ch[x][!d]) fa[ch[x][!d]]=p;
    ch[x][!d]=p; fa[p]=x;
}
void splay(int x){
    while(!isRoot(x)){ int p=fa[x], g=fa[p];
        if(isRoot(p)) rot(x);
        else if((ch[g][0]==p)==(ch[p][0]==x)){ rot(p); rot(x); }
        else { rot(x); rot(x); } }
}
void access(int x){                       // make root->x a single preferred path
    for(int last=0; x; last=x, x=fa[x]){  // x climbs the path-parent spine
        splay(x);
        ch[x][1]=last;                    // deeper segment becomes non-preferred
    }
}
int findroot(int x){                       // represented root = shallowest on x's path
    access(x);
    while(ch[x][0]) x=ch[x][0];            // leftmost in the aux = min depth
    splay(x); return x;
}
\end{lstlisting}

\begin{figure}[H]
\centering
\begin{tikzpicture}[diagram, scale=0.8, transform shape,
   snode/.style={circle, draw=primaryblue!70, fill=primaryblue!8, line width=0.7pt,
                 minimum size=6mm, font=\scriptsize\bfseries, inner sep=1pt},
   pref/.style={circle, draw=difforange, fill=difforange!16, line width=1pt,
              minimum size=6.4mm, font=\footnotesize\bfseries, inner sep=1pt}]
  % represented tree R
  \node[pref] (n1) at (0,0) {1};
  \node[pref] (n2) at (-2,-1.4) {2};
  \node[snode] (n3) at (2.5,-1.4) {3};
  \node[snode] (n4) at (-3.2,-2.8) {4};
  \node[pref] (n5) at (-0.9,-2.8) {5};
  \node[snode] (n8) at (2.5,-2.8) {8};
  \node[snode] (n6) at (-1.8,-4.2) {6};
  \node[pref] (n7) at (0.0,-4.2) {7};
  \draw[difforange, line width=1.6pt] (n1)--(n2) (n2)--(n5) (n5)--(n7);
  \draw (n1)--(n3) (n2)--(n4) (n3)--(n8) (n5)--(n6);
  \node[annot, font=\scriptsize, text=difforange, anchor=west] at (1.4,0.4) {one aux splay (depth order): 1,2,5,7};
  \node[annot, font=\scriptsize, anchor=west] at (3.4,-1.4) {path-parent 3 \(\to\) 1};
  \node[annot, font=\scriptsize, anchor=west] at (-4.9,-2.8) {4 \(\to\) 2};
  \node[annot, font=\scriptsize, anchor=west] at (1.4,-4.2) {6 \(\to\) 5};
\end{tikzpicture}
\caption{The running forest \(R\) of \(\figref{c23-tree}\) after \code{access(7)}.
The solid orange preferred path \(1\!-\!2\!-\!5\!-\!7\) is a single auxiliary splay
(aux nodes in orange, ordered by depth). The three off-path subtrees \(\{3,8\}\),
\(\{4\}\), \(\{6\}\) are their own preferred paths; each path's top node keeps a
dashed \emph{path-parent} pointer into the spine (\(3\!\to\!1\), \(4\!\to\!2\),
\(6\!\to\!5\)). These dashed pointers are \(\mathit{fa}\) links that are
\emph{not} aux children -- that is exactly what \code{is\_root} tests.}
\label{fig:c23-aux}
\end{figure}

\begin{theorem}{Two accesses read the lowest common ancestor}{c23-access-lca}
In a rooted link-cut tree (no \code{makeroot}), for nodes \(u,v\) in the same
represented tree, after \code{access(u)} a path-parent climb from \(v\)
\(\big(\)for \(x{=}v\); \(x\neq 0\); \(x{=}\mathit{fa}[x]\) after \code{splay(x)}\(\big)\)
has its \emph{last} visited node equal to \(\lca(u,v)\).
\end{theorem}
\begin{proofbox}
\code{access(u)} turns the whole root-to-\(u\) path into one aux tree whose top
(the represented root) has \(\mathit{fa}=0\). The climb from \(v\) visits, in
bottom-to-top order, one node from each preferred path that \(v\)'s ancestor chain
passes through, because after \code{splay(x)} the value \(\mathit{fa}[x]\) is the
path-parent, i.e.\ the real-tree parent of \(x\)'s path top -- never skipping an
ancestor. Every visited node is an ancestor of \(v\). The loop stops the first time
the visited node lies inside the root-to-\(u\) aux tree, since splaying such a node
makes it that aux tree's root and forces \(\mathit{fa}=0\). That node is the
deepest ancestor of \(v\) that also lies on root-to-\(u\), i.e.\ an ancestor of
\(u\); being reachable from \(v\)'s chain and from \(u\)'s chain and maximal with
that property, it is precisely \(\lca(u,v)\). \(\Box\)
\end{proofbox}

\begin{complexitybox}[Access and LCA]
\textbf{Per operation} \(\bigO{\log n}\) amortised: an \code{access} is a sequence
of splays along the spine, and (Sleator--Tarjan) the total splaying cost is
logarithmic amortised -- the same potential as \(\thmref{c23-splay-amort}\) applied to
each aux tree. \code{findroot} is one \code{access} plus a left-descent. The LCA
query is \code{access(u)} plus a climb that is itself an \code{access}-style spine,
so also \(\bigO{\log n}\) amortised. \textbf{Space} \(\bigO{n}\) for
\(\mathit{ch},\mathit{fa}\).
\end{complexitybox}

% ---------------------------------------------------------------
\subsection{Problem Lab: SPOJ DYNALCA -- LCA while the forest is rewired}
% ---------------------------------------------------------------

\solvedproblem{c23-dynalca}{Dynamic LCA}{SPOJ DYNALCA}
  {https://www.spoj.com/problems/DYNALCA/}{SPOJ classic L5}
  {rooted link-cut tree; access; findroot; LCA by a path-parent climb}

\begin{problemstatement}[Dynamic LCA (statement and sample verified against the live SPOJ page)]
A forest starts as \(N\) single-vertex rooted trees (\(1\le N\le 10^5\)), vertices
\(1..N\). Process \(M\le 10^5\) queries:
\begin{itemize}[nosep]
  \item \texttt{link A B}: add an edge making \(A\) a child of \(B\); \(A\) is
        guaranteed to be a root, and \(A,B\) are in different trees.
  \item \texttt{cut A}: remove the edge from \(A\) to its parent; \(A\) is
        guaranteed non-root. Its subtree detaches into a new tree.
  \item \texttt{lca A B}: print the lowest common ancestor of \(A\) and \(B\), which
        are guaranteed to be in the same tree.
\end{itemize}

\textbf{Example (official).}
\begin{verbatim}
Input:              Output:
5 9                 1
lca 1 1             2
link 1 2            3
link 3 2            2
link 4 3
lca 1 4
lca 3 4
cut 4
link 5 3
lca 1 5
\end{verbatim}
\end{problemstatement}

\begin{attemptgate}[Attempt it yourself first -- \lv{5}, give yourself 35 minutes]
A rooted LCT has the pieces already: \code{access}, \code{findroot}, and the
\(\thmref{c23-access-lca}\) climb. \peek{\texttt{link A B} with \(A\) a root is just
\code{access(A); fa[A]=B;}.}
\nudge{\texttt{cut A}: \code{access(A)} makes root\(\to\)A one path with A the aux
root; A's left child is the whole ancestor block, so detach it
(\(\mathit{ls}[A]=0,\ \mathit{fa}[\text{old-}ls]=0\)).}
\fullhint{For \texttt{lca A B}: run \code{access(A)}; then climb from \(B\) with
\code{for(x=B; x; x=fa[x]) splay(x), last=x;}. The recorded \texttt{last} is
\(\lca(A,B)\) (\(\thmref{c23-access-lca}\)). Because DYNALCA never reroots, there is
no reverse lazy tag and no \code{pushdown} to worry about -- the cleanest possible
LCT.}
\end{attemptgate}

\solstep{Understanding}
Every query is a rooted-LCT primitive: link attaches a root's path-parent, cut
detaches a node from its ancestor block, and lca is the two-access trick of
\(\thmref{c23-access-lca}\). No \code{makeroot}, no aggregates, no lazy propagation
-- the forest simply stays rooted and the aux splays track preferred paths. This is
DYNALCA's gift to a first-time LCT author: it exercises the whole \code{access}
machinery with none of the evert subtlety that makes general LCTs error-prone.

\solstep{Key Observations}
\begin{observation}
\texttt{link A B} with \(A\) a represented root: after \code{access(A)}, \(A\) is an
aux root with \(\mathit{fa}[A]=0\); writing \(\mathit{fa}[A]=B\) turns the singleton
path \(A\) into a child whose path-parent is \(B\). No restructuring of \(B\) is
needed because a path-parent may aim at any node, preferred or not.
\end{observation}
\begin{observation}
\texttt{cut A} (\(A\) non-root): \code{access(A)} leaves \(A\) as aux root of the
path root\(\to A\), with its left child being exactly the ancestor block
root\(\to\!\text{parent}(A)\). Clearing \(\mathit{ls}[A]\) and its back-pointer severs
the single represented edge \((\text{parent}(A),A)\) and strands \(A\)'s (empty, here)
subtree as a new root -- the ancestors stay connected to each other inside the
detached block.
\end{observation}

\solstep{The Approach Ladder}
\begin{center}\renewcommand{\arraystretch}{1.25}\setlength{\tabcolsep}{3pt}
\begin{tabular}{@{}c p{0.26\linewidth} p{0.30\linewidth} l c@{}}
\toprule
\textbf{\#} & \textbf{Approach} & \textbf{Idea} & \textbf{Time} & \textbf{Verdict}\\
\midrule
1 & \lv{3}\ BFS per lca & climb parents to root by brute force & \(O(N)\)/query & \tle \\
2 & \lv{4}\ binary lifting & \(\lca\) tables (\chref{c11}) & \(O(\log n)\)/query & \textcolor{warningred}{\ding{55}\ \textbf{stale}} \\
3 & \lv{5}\ \textbf{rooted LCT} & access + path-parent climb (\(\coderef{c23-lct-core}\)) & \(\bigO{\log n}\) amort & \ac \\
\bottomrule
\end{tabular}\end{center}
Rung~2 is the tempting wrong tool: binary lifting answers a static LCA fast, but
a single \texttt{link}/\texttt{cut} invalidates every descendant's table and forces
an \(O(n\log n)\) rebuild -- hopeless with \(10^5\) edits. Rung~3 pays the topology
change in the same \(O(\log n)\) as the query.

\solstep{Solution}
Maintain the \(\coderef{c23-lct-core}\) arrays. Map:
\texttt{link A B} \(\to\) \code{access(A); fa[A]=B};
\texttt{cut A} \(\to\) \code{access(A); if(ls) detach};
\texttt{lca A B} \(\to\) \code{access(A); climb from B, return last}.

\solstep{Implementation}
\begin{lstlisting}[caption={SPOJ DYNALCA: the rooted LCT of lst:c23-lct-core plus link/cut/lca. Reuses isRoot/rot/splay/access/findroot verbatim. ASCII only}, label={lst:c23-dynalca}]
#include <bits/stdc++.h>
using namespace std;
const int MAXN = 100005;
int ch[MAXN][2], fa[MAXN];
#define ls(x) ch[x][0]
#define rs(x) ch[x][1]
// ... paste isRoot, rot, splay, access, findroot from lst:c23-lct-core ...
void link(int a, int b){ access(a); fa[a]=b; }          // a is a represented root
void cut(int a){ access(a); if(ls(a)){ fa[ls(a)]=0; ls(a)=0; } } // sever a->parent
int lca(int a, int b){
    if(findroot(a)!=findroot(b)) return 0;              // (guaranteed same tree here)
    access(a);
    int last=0;
    for(int x=b; x; x=fa[x]){ splay(x); last=x; }       // Thm: last == lca(a,b)
    splay(last);
    return last;
}
int main(){
    int N,M;  scanf("%d %d",&N,&M);
    char op[8];
    while(M--){
        scanf("%s",op);
        if(op[0]=='l' && op[1]=='i'){ int A,B; scanf("%d %d",&A,&B); link(A,B); }
        else if(op[0]=='c'){ int A; scanf("%d",&A); cut(A); }
        else { int A,B; scanf("%d %d",&A,&B); printf("%d\n", lca(A,B)); }
    }
}
\end{lstlisting}

\solstep{Trace}
\begin{tracebox}[Official DYNALCA sample, rooted-LCT state]
Start: five singletons \(1..5\), all roots.\\
\(\texttt{lca }1\ 1\): \code{findroot(1)}=\(1\) both sides; \code{access(1)}; climb
from \(1\) reaches \(1\) (\(\mathit{fa}=0\)) \(\Rightarrow\) \(\boxed{1}\).\\
\(\texttt{link }1\ 2\): \(\mathit{fa}[1]{=}2\). \(\texttt{link }3\ 2\):
\(\mathit{fa}[3]{=}2\). \(\texttt{link }4\ 3\): \(\mathit{fa}[4]{=}3\). Tree rooted at
\(2\): \(2\!\to\!\{1,3\}\), \(3\!\to\!4\).\\
\(\texttt{lca }1\ 4\): \code{access(1)} makes path \(2\!-\!1\) preferred
(\(\mathit{fa}\) chain \(1\!\to\!2\)). Climb from \(4\): \(splay(4)\) (its own path),
\(\mathit{fa}[4]{=}3\) \(\to splay(3)\), \(\mathit{fa}[3]{=}2\) \(\to splay(2)\): \(2\)
is in the root\(\to1\) aux, so \(\mathit{fa}[2]{=}0\), loop ends,
\(\mathit{last}=2\) \(\Rightarrow\) \(\boxed{2}\).\\
\(\texttt{lca }3\ 4\): \code{access(3)} (path \(2\!-\!3\)). Climb from \(4\):
\(splay(4)\), \(\mathit{fa}[4]{=}3\) \(\to splay(3)\) joins the aux,
\(\mathit{fa}[3]{=}0\) \(\Rightarrow\) \(\mathit{last}=3\), \(\boxed{3}\).\\
\(\texttt{cut }4\): \code{access(4)} (path \(2\!-\!3\!-\!4\)), clear
\(\mathit{ls}[4]\) \(\Rightarrow\) \(4\) becomes a root, \(2\!-\!3\) intact.\\
\(\texttt{link }5\ 3\): \(\mathit{fa}[5]{=}3\).\\
\(\texttt{lca }1\ 5\): \code{access(1)} (path \(2\!-\!1\)). Climb from \(5\):
\(splay(5)\), \(\mathit{fa}[5]{=}3\) \(\to splay(3)\), \(\mathit{fa}[3]{=}2\) \(\to
splay(2)\) joins aux, \(\mathit{fa}[2]{=}0\) \(\Rightarrow\) \(\boxed{2}\).\\
Printed: \(1,2,3,2\) -- matching the official output.
\end{tracebox}

\begin{warningbox}[access is not ``splay to the represented root'']
The loop variable in \(\coderef{c23-lct-core}\)'s \code{access} climbs
\(\mathit{fa}\) \emph{after} each splay, so it hops across path-parent edges -- but
the very first \code{splay(x)} makes \(x\) the aux root, so \(\mathit{fa}[x]\) at
that moment is the path-parent, not a splay parent. Reading \(\mathit{fa}[x]\)
\emph{before} splaying, or forgetting to re-read it after, walks a wrong spine and
silently corrupts the preferred-path split (WA that depends on access order). Also,
in \texttt{cut} you must \code{access} first; cutting a raw edge pointer leaves a
dangling aux child and breaks \code{is\_root}.
\end{warningbox}

\begin{hardeningbox}
\begin{itemize}[nosep,leftmargin=1.4em]
  \item \textbf{No lazy tag here.} DYNALCA is rooted; do NOT introduce
        \code{makeroot}/reverse, which would demand \code{pushdown} inside every
        \code{splay} before rotating. Adding evert ``because LCT usually has it'' is
        the classic way to break an otherwise-correct rooted LCT.
  \item \textbf{findroot before lca.} Even though DYNALCA promises same-tree
        queries, keeping \(\texttt{if(findroot(a)!=findroot(b)) return 0}\) is a cheap
        guard that also makes the template reusable for connectivity problems.
  \item \textbf{access(a) then access(b) is not the LCA.} A common mis-mnemonic:
        two plain \code{access} calls leave \(b\) on top, not the lca. Use the
        \emph{climb-and-record} of \(\thmref{c23-access-lca}\), or switch to the
        \code{makeroot}-free \code{lca} in \(\coderef{c23-dynalca}\) verbatim.
\end{itemize}
\end{hardeningbox}

\begin{recallbox}[25.2 in one breath]
\emph{Store the rooted forest as one aux splay per preferred path (ordered by
depth); \(\mathit{fa}\) is an aux-parent or a path-parent, told apart by
\code{is\_root}. \code{access(v)} walks the path-parent spine, splaying and
reattaching right children, until root\(\to\)v is one contiguous splay. Then
\code{findroot} is a left-descent, and \(\lca(u,v)\) is \code{access(u)} plus a
path-parent climb from \(v\) whose last splayed node is the answer
(\(\thmref{c23-access-lca}\)). SPOJ~DYNALCA is exactly link/cut/lca in that rooted
setting -- the splay engine of 25.1, now driving a changing forest.}
\selftest{
  \item What does \code{is\_root(x)} distinguish? \emph{(aux-parent vs path-parent in the shared \(\mathit{fa}\) pointer, Def.~\defref{c23-prefer}.)}
  \item After \code{access(v)}, what does \(v\)'s aux tree contain? \emph{(the whole root-to-\(v\) path, in depth order, \(v\) splayed to the top.)}
  \item Give the two steps that read \(\lca(u,v)\) with no rerooting. \emph{(\code{access(u)}; climb \(v\)'s path-parents, last splayed node = lca, \(\thmref{c23-access-lca}\).)}
}
\end{recallbox}

% =====================================================================
%  25.3  Evert, Link, Cut, and Path Aggregates -- the full LCT.
%  Exports: sec:c23-makeroot  def:c23-makeroot  def:c23-evert-agg
%           thm:c23-path-query  lst:c23-evert  lst:c23-pathsum
%  WORKED example only -- NO solvedproblem (the in-section solved count
%  stays 2; the flagship dynamic-path problem is the chapter Summit).
%  The worked task is a self-contained "dynamic tree vertex add / path sum":
%  link, cut, add-to-vertex, query sum along a path, fully traced on a
%  4-vertex chain so the reader sees makeroot expose a path and the reverse
%  lazy tag do its work. Modelled on the standard rooted-then-everted LCT.
% =====================================================================
\section{Evert, Link, Cut, and Path Aggregates}\label{sec:c23-makeroot}

\begin{previouslybox}[Ch.~22 needed a fixed root; the LCT gets to pick one per query]
In \(\secref{c23-prefer}\) the forest stayed rooted and \code{access} could only
expose a path \emph{downward} from that fixed root. To link two arbitrary trees,
cut a single edge, or query the path between two nodes that are neither an
ancestor of the other, we need one more move: \emph{reroot the represented tree at
any node}. Chapter~22 sidestepped this by choosing a root once in its DFS
(\(\secref{c16-decomp}\)); the LCT does it on demand with \code{makeroot}, and the
price is a single lazy flag. This section adds \code{makeroot}, the reverse tag,
\code{link}/\code{cut}, and a path aggregate -- completing the tool the Summit
will ride on a block-cut tree.
\end{previouslybox}

\begin{topicintro}
\code{makeroot(v)} = \code{access(v)} followed by flipping a \emph{reverse} lazy tag
on \(v\)'s aux tree. Since that aux tree currently stores the depth-ordered
root-to-\(v\) path, reversing it makes the old root the deepest node and \(v\) the
new root -- inverting the whole path's parent relations at once, exactly like
reversing a sequence in an implicit splay. Once you can evert, the three composite
operations are one-liners: \code{link(u,v)} everts \(u\) then hangs it under \(v\);
\code{cut(u,v)} everts \(u\), exposes \(v\), and checks the two are adjacent; a path
aggregate \code{query(u,v)} everts \(u\), accesses \(v\), and reads \(v\)'s splay
fold. The fold is any monoid (sum, min, max, gcd) carried up the aux trees by
\code{pull}, and pushed down through the reverse tag by \code{pushdown}.

\textbf{Prerequisites.} \(\secref{c23-prefer}\) -- \code{access}, \code{splay},
\code{is\_root}, \code{findroot}. \(\secref{c23-splay}\) -- lazy-tag-free splaying,
now augmented. The monoid/segment-fold idea is Ch.~11 (\chref{c14}); nothing about it
is re-taught.

\textbf{After this section you can}
\begin{itemize}[nosep]
  \item write \code{makeroot} as \code{access} + a reverse tag, and explain why
        reversal reroots; \hfill\lv{5}
  \item implement \code{link}/\code{cut}/\code{query} from \code{makeroot}; \hfill\lv{5}
  \item carry a vertex-sum (or min) aggregate through rotations and lazy flips; \hfill\lv{5}
  \item trace a dynamic path-sum: evert, expose, fold, read. \hfill\lv{5}
\end{itemize}
\end{topicintro}

\begin{definition}{Evert and the reverse lazy tag}{c23-makeroot}
\code{makeroot(v)} performs \code{access(v)} (so the root-to-\(v\) path is one aux
splay with \(v\) at its root and depth increasing left-to-right) and then applies a
\term{reverse} tag to \(v\)'s splay: swap its two children and mark the tag to be
pushed down. Inverting the in-order sequence inverts the path's direction, so the
old root becomes a leaf and \(v\) becomes the represented root. The tag is a plain
lazy propagation -- identical to reversing a segment in an implicit splay.
\end{definition}

\begin{lstlisting}[caption={Adding evert power to the LCT core: pull, apply-reverse, pushdown, makeroot, link, cut, and a path-sum query. Children ls/rs, aggregate sum[], per-node val[], lazy rev[]. Call pushdown top-down inside splay before rotating. ASCII only}, label={lst:c23-evert}]
int ch[MAXN][2], fa[MAXN], rev[MAXN];
long long val[MAXN], sum[MAXN];
void pull(int x){ sum[x] = val[x] + sum[ch[x][0]] + sum[ch[x][1]]; }
void applyRev(int x){ if(!x) return; swap(ch[x][0], ch[x][1]); rev[x]^=1; }
void pushdown(int x){ if(rev[x]){ applyRev(ch[x][0]); applyRev(ch[x][1]); rev[x]=0; } }
// splay() from lst:c23-lct-core, but push the whole ancestor chain first:
void pushPath(int x){ if(!isRoot(x)) pushPath(fa[x]); pushdown(x); }
//   ... then in splay(x): call pushPath(x) before each rot(...) group.
void makeroot(int x){ access(x); applyRev(x); }
int findroot(int x){ access(x); while(ch[x][0]){ pushdown(x); x=ch[x][0]; } splay(x); return x; }
void link(int u,int v){ makeroot(u); if(findroot(v)!=u) fa[u]=v; }
void cut(int u,int v){
    makeroot(u); access(v);            // path u->v is v's aux tree, u is leftmost
    if(ch[v][0]==u && !ch[u][1]){ ch[v][0]=0; fa[u]=0; pull(v); }
}
long long query(int u,int v){          // sum of val on the simple path u..v
    makeroot(u); access(v);            // u..v now a single aux tree rooted at v
    return sum[v];
}
void addVal(int u,long long d){ access(u); val[u]+=d; pull(u); }
\end{lstlisting}

\begin{theorem}{Evert-then-access exposes a path, and the fold answers it}{c23-path-query}
After \code{query(u,v)} = \code{makeroot(u)} then \code{access(v)}, the auxiliary
splay tree rooted at \(v\) contains exactly the vertices of the simple path
\(u\leadsto v\) in the represented tree, and \(\mathit{sum}[v]\) equals the sum of
their values.
\end{theorem}
\begin{proofbox}
\code{makeroot(u)} makes \(u\) the represented root (\(\defref{c23-makeroot}\)), so
\(v\)'s ancestors run up to \(u\). \code{access(v)} then makes the current
root-to-\(v\) path -- which is precisely \(u\leadsto v\) -- a single preferred path,
i.e.\ one aux splay with \(v\) splayed to the top (\(\defref{c23-access}\)); no
other vertex is in that splay, because access detaches every right child along the
spine, leaving only the path. The in-order sequence of the splay is the path by
depth, and \code{pull} maintains \(\mathit{sum}[x]=\mathit{val}[x]+
\mathit{sum}[\mathit{ls}]+\mathit{sum}[\mathit{rs}]\) through every rotation and
lazy reversal, so the aggregate of the root \(v\) is the sum over the whole splay,
i.e.\ over the path. \(\Box\)
\end{proofbox}

\begin{complexitybox}[Evert, link, cut, query]
\textbf{Per operation} \(\bigO{\log n}\) amortised -- \code{makeroot} is one
\code{access} plus an \(O(1)\) tag; \code{link}/\code{cut}/\code{query} are
constant many accesses/splays. \textbf{Space} \(\bigO{n}\). The reverse tag costs
nothing asymptotically: it is one swap at a node, propagated lazily, and any
aggregate that is invariant under reordering (sum, min, max, gcd, xor) needs no
order-aware fix-up -- only an order-\emph{sensitive} fold would have to also flip
its combine direction inside \code{applyRev}.
\end{complexitybox}

% ---------------------------------------------------------------
\subsection{Worked Example: dynamic vertex-add / path-sum}
% ---------------------------------------------------------------

\begin{problemstatement}[A self-contained dynamic-tree exercise (worked, not a judge problem)]
A forest of weighted vertices starts as isolated nodes. Support online:
\(\texttt{link}(u,v)\) join two trees by an edge; \(\texttt{cut}(u,v)\) remove an
edge; \(\texttt{add}(u,d)\) add \(d\) to vertex \(u\)'s value; \(\texttt{query}(u,v)\)
print the sum of values along the unique path \(u\leadsto v\). All in
\(O(\log n)\) amortised.
\end{problemstatement}

\begin{insightbox}[The whole program is \code{makeroot} + \code{access}]
Add is \code{access(u)} then bump \(\mathit{val}[u]\). Link, cut and query all start
by \code{makeroot}-ing one endpoint so the path becomes root-to-the-other and a
final \code{access} folds it. Once you internalise ``\code{makeroot(u); access(v)}
turns \(u\leadsto v\) into \(v\)'s splay'', every dynamic path aggregate is the same
two lines -- only the \code{pull} monoid changes (sum here, min at the Summit).
\end{insightbox}

\begin{tracebox}[Build a path \(1\!-\!2\!-\!3\!-\!4\), then query and reroot]
Values \(v_1{=}1,v_2{=}2,v_3{=}3,v_4{=}4\).\\
\(\texttt{link}(1,2)\): \code{makeroot(1)} (\(1\) alone, tag on a singleton is a
no-op after pushdown), \(\mathit{fa}[1]{=}2\).\\
\(\texttt{link}(2,3)\): \code{makeroot(2)}; \(2\)'s path is \(2\!-\!1\) reversed to
root \(2\); hang \(2\) under \(3\): \(\mathit{fa}[2]{=}3\). Tree now \(3\!-\!2\!-\!1\).\\
\(\texttt{link}(3,4)\): \code{makeroot(3)} roots the chain at \(3\); \(\mathit{fa}[3]{=}4\).
Represented chain \(4\!-\!3\!-\!2\!-\!1\).\\
\(\texttt{query}(1,4)\): \code{makeroot(1)} re-roots so \(1\) is root
(\(1\!-\!2\!-\!3\!-\!4\)); \code{access(4)} exposes path \(1\!-\!2\!-\!3\!-\!4\) as
\(4\)'s splay; \(\mathit{sum}[4]=1{+}2{+}3{+}4=\boxed{10}\).\\
\(\texttt{add}(3,5)\): \code{access(3)}; \(\mathit{val}[3]=8\); ancestors' sums fixed.\\
\(\texttt{query}(1,4)\) again: path sum \(1{+}2{+}8{+}4=\boxed{15}\).\\
\(\texttt{cut}(2,3)\): \code{makeroot(2)}; \code{access(3)}; check \(\mathit{ch}[3][0]{=}2\)
and \(2\) has no right child; detach \(\Rightarrow\) trees \(\{1,2\}\) and \(\{3,4\}\).\\
\(\texttt{query}(1,2)\): \(\boxed{3}\) (still connected); \(\texttt{query}(1,4)\) now
returns nothing useful -- guard with \(\mathit{findroot}(1)\ne\mathit{findroot}(4)\).
\end{tracebox}

\begin{lstlisting}[caption={The worked path-sum driver (main only; reuse lst:c23-evert verbatim). Queries mirror YS ``dynamic tree vertex add path sum''. ASCII only}, label={lst:c23-pathsum}]
int main(){
    int n,q;  scanf("%d %d",&n,&q);
    for(int i=1;i<=n;i++){ val[i]=/* read */; sum[i]=val[i]; }
    while(q--){
        char op[8]; scanf("%s",op);
        if(op[0]=='l'){ int u,v; scanf("%d %d",&u,&v); link(u,v); }
        else if(op[0]=='c'){ int u,v; scanf("%d %d",&u,&v); cut(u,v); }
        else if(op[0]=='a'){ int u; long long d; scanf("%d %lld",&u,&d); addVal(u,d); }
        else { int u,v; scanf("%d %d",&u,&v);
               if(findroot(u)!=findroot(v)) printf("disconnected\n");
               else printf("%lld\n", query(u,v)); }
    }
}
\end{lstlisting}

\begin{warningbox}[Push the tags BEFORE you rotate -- every LCT bug starts here]
\code{pushPath(x)} must run \emph{top-down} over the ancestor chain before
\code{splay(x)}'s rotations, or a pending reverse tag sits above \(x\) and the
rotation reads a stale child pointer -- the tree silently loses or duplicates path
segments (wrong answer only on paths that cross an everted region, so it hides for
dozens of operations). Pair it with the mirror discipline in \code{makeroot}: flip
\(\mathit{rev}\) and swap children \emph{together}, and never call \code{pull}
between the swap and a needed \code{pushdown}.
\end{warningbox}

\begin{hardeningbox}
\begin{itemize}[nosep,leftmargin=1.4em]
  \item \textbf{Guard link/cut.} \code{link} should skip when \(u,v\) are already
        connected (\(\mathit{findroot}(v)==u\)) or you create a cycle and corrupt the
        structure; \code{cut} must verify \(\mathit{ch}[v][0]==u\) and \(\mathit{ch}[u][1]==0\)
        before detaching, so a non-edge cut is a harmless no-op.
  \item \textbf{Order-sensitive folds need a second aggregate.} Sum/min/max/gcd are
        commutative, so one \code{sum[]} suffices. For a fold like
        ``string along the path'' you must store both forward and reversed folds and
        swap them in \code{applyRev}.
  \item \textbf{findroot's pushdown.} The left-descent in \code{findroot} has to
        \code{pushdown} at each step (see \(\coderef{c23-evert}\)); a pending reversal
        makes the ``leftmost'' pointer wrong and returns the wrong root.
\end{itemize}
\end{hardeningbox}

\begin{recallbox}[25.3 in one breath]
\emph{\code{makeroot(u)} = \code{access(u)} + a reverse tag, turning \(u\) into the
represented root. Then \code{link}/\code{cut}/\code{query} are ``evert one endpoint,
access the other'': \code{query(u,v)} folds \(v\)'s aux splay, which after
\code{makeroot(u); access(v)} is exactly the path \(u\leadsto v\)
(\(\thmref{c23-path-query}\)). Push lazy tags top-down before rotating or the structure
rots out. This full LCT is what the chapter Summit mounts on a block-cut tree.}
\selftest{
  \item What two calls expose path \(u\leadsto v\) as one aux splay? \emph{(\code{makeroot(u)} then \code{access(v)}, \(\thmref{c23-path-query}\).)}
  \item Why is a reverse tag enough to reroot? \emph{(the aux splay is the depth-ordered path; inverting the sequence inverts parent direction, Def.~\defref{c23-makeroot}.)}
  \item Which aggregates are free to fold, and which need extra care? \emph{(commutative ones -- sum/min/max/gcd/xor -- are free; order-sensitive folds need a reversed twin, Hardening.)}
}
\end{recallbox}

% =====================================================================
%  25.4  Beyond Paths -- virtual children, subtree aggregates, and the
%  connectivity question (LCT vs Euler-tour tree vs DSU).
%  Exports: sec:c23-beyond  def:c23-virtual  def:c23-map
%           thm:c23-vsum  lst:c23-subtree  fig:c23-virtual
%  WORKED examples only -- NO solvedproblem (in-section solved stays 2).
% =====================================================================
\section{Beyond Paths: Virtual Children and Connectivity}\label{sec:c23-beyond}

\begin{previouslybox}[Paths were easy because a splay only sees the solid chain]
\(\secref{c23-makeroot}\)'s \code{query} read a path and was done, because a
preferred path \emph{is} the splay subtree -- the off-path branches simply did not
matter. The moment you ask for a \emph{subtree} aggregate on a moving tree, those
off-path branches matter enormously: an \code{access} is constantly promoting one
child to solid and demoting the others to \emph{virtual}, and a node's splay
children are only ever its preferred neighbours. This section shows the fix -- a
separate accumulator \(\mathit{vir}[\cdot]\) for virtual children, kept current on
every access -- and then steps back to ask, honestly, when you should not be
reaching for an LCT at all.
\end{previouslybox}

\begin{topicintro}
Two limits define the LCT's real job description. First, a subtree aggregate: a
node's aux splay holds only its \emph{preferred} child, so the other represented
children are ``virtual'' and invisible to \code{pull}; you fix this by storing their
combined value in \(\mathit{vir}[x]\) and moving a child between the solid and
virtual sets during \code{access} (\(\thmref{c23-vsum}\)). Second, connectivity: if the
query only asks ``are \(u,v\) in the same tree?'' -- never a path value -- an
Euler-tour tree or a (rollback) DSU is smaller, faster, and far less bug-prone.
Knowing both limits is what separates an LCT that ships from an LCT that TLEs on a
\(10^5\) connectivity problem it was never meant to solve.

\textbf{Prerequisites.} \(\secref{c23-makeroot}\) -- \code{access}/\code{makeroot}
and a \code{pull} monoid. Chapter~5 (\chref{c24}) -- the Euler tour that an
Euler-tour tree keeps in a balanced BST. Chapter~17 (\chref{c28}) -- DSU, and
Chapter~24 (\chref{c18}) -- the offline rollback / divide-and-conquer alternative.

\textbf{After this section you can}
\begin{itemize}[nosep]
  \item keep a virtual-child accumulator \(\mathit{vir}\) correct across accesses; \hfill\lv{5}
  \item answer a subtree-sum on a link/cut tree; \hfill\lv{5}
  \item decide LCT vs Euler-tour tree vs DSU from the query's shape. \hfill\lv{5}
\end{itemize}
\end{topicintro}

\begin{definition}{Virtual children and the \(\mathit{vir}\) accumulator}{c23-virtual}
Call a represented child of \(x\) \term{solid} if it lies on \(x\)'s current
preferred path (so it is one of \(x\)'s two aux children); every other represented
child is \term{virtual}. Define \(\mathit{vir}[x]\) as the fold (here a sum) over the
\emph{subtree-aggregate} of each virtual child. A node's full subtree value is then
\(\mathit{sub}[x]=\mathit{val}[x]+\mathit{vir}[x]+\mathit{sub}[\text{solid children}]\),
and the path fold \(\mathit{sum}[x]\) of \(\secref{c23-makeroot}\) is the special case
that omits \(\mathit{vir}\).
\end{definition}

\begin{theorem}{Access moves children between the solid and virtual sets}{c23-vsum}
In \code{access}, when at node \(x\) we set the right child from an old value \(r\) to
the new preferred child \(y\), \(r\) becomes virtual and \(y\) stops being virtual.
Maintaining
\[
  \mathit{vir}[x]\;\mathrel{-}=\;\mathit{sub}[y]\quad(\text{if }y\ne 0),\qquad
  \mathit{vir}[x]\;\mathrel{+}=\;\mathit{sub}[r]\quad(\text{if }r\ne 0)
\]
at that instant, and folding \(\mathit{vir}\) into \(\mathit{sub}\) (not into the pure
path \(\mathit{sum}\)) keeps every subtree aggregate exact after any access.
\end{theorem}
\begin{proofbox}
A subtree aggregate must count every represented child's subtree exactly once. At
any moment the represented children of \(x\) partition into the (at most one) solid
child on the preferred path and the rest, which are virtual. Before the access step,
\(r\) was the solid (right/deeper) child and \(y\) was virtual; after, \(y\) is solid
and \(r\) is virtual. Removing \(\mathit{sub}[y]\) from \(\mathit{vir}[x]\) (it is now
counted via the solid splay path instead) and adding \(\mathit{sub}[r]\) (it is no
longer on the path) preserves the invariant ``\(\mathit{vir}[x]\) is the sum over
precisely the virtual children's subtrees.'' Since \(\mathit{sub}[x]\) is defined as
\(\mathit{val}+\mathit{vir}+\) solid-subtrees, and every rotation of the aux splay only
reshuffles solid children (whose \(\mathit{sub}\) values are unchanged), \(\mathit{sub}\)
stays exact. \(\Box\)
\end{proofbox}

\begin{figure}[H]
\centering
\begin{tikzpicture}[diagram, scale=0.82, transform shape,
   snode/.style={circle, draw=primaryblue!70, fill=primaryblue!8, line width=0.7pt,
                 minimum size=6mm, font=\scriptsize\bfseries, inner sep=1pt},
   virt/.style={circle, draw=successgreen!75, fill=successgreen!10, line width=1pt,
              minimum size=6.4mm, font=\footnotesize\bfseries, inner sep=1pt},
   pref/.style={circle, draw=difforange, fill=difforange!16, line width=1pt,
              minimum size=6.4mm, font=\footnotesize\bfseries, inner sep=1pt}]
  \node[pref] (n2) at (0,0) {2};
  \node[pref] (n5) at (0,-1.5) {5};
  \node[virt] (n4) at (-2.2,-1.5) {4};
  \node[virt] (n6) at (2.2,-1.5) {6};
  \node[pref] (n7) at (0,-3.0) {7};
  \draw[difforange, line width=1.6pt] (n2)--(n5)--(n7);
  \draw[successgreen, line width=1.3pt, densely dashed] (n2)--(n4);
  \draw[successgreen, line width=1.3pt, densely dashed] (n5)--(n6);
  \node[annot, font=\scriptsize, text=successgreen, anchor=east] at (-2.4,-0.7) {virtual: in \(\mathit{vir}[2]\)};
  \node[annot, font=\scriptsize, text=successgreen, anchor=west] at (2.4,-0.7) {virtual: in \(\mathit{vir}[5]\)};
\end{tikzpicture}
\caption{On the running forest \(R\) (\(\figref{c23-tree}\)) after \code{access(7)},
the solid path \(2\!-\!5\!-\!7\) (orange) lives in the aux splay; children \(4\) (of
\(2\)) and \(6\) (of \(5\)) are \emph{virtual} (dashed green) and are excluded from
the path fold but included in \(\mathit{vir}[2]\), \(\mathit{vir}[5]\) so that a
\emph{subtree} query still counts them. When a later access demotes \(5\)'s solid
child and promotes \(6\), the two \(\mathit{sub}\) values trade between the splay and
\(\mathit{vir}[5]\) exactly as \(\thmref{c23-vsum}\) prescribes.}
\label{fig:c23-virtual}
\end{figure}

\begin{lstlisting}[caption={Subtree-sum-capable LCT: path fold sum[] (solid only) and subtree fold sub[] (solid + virtual) kept side by side. access() now updates vir[] as children swap between solid and virtual. ASCII only}, label={lst:c23-subtree}]
long long val[MAXN], sum[MAXN], vir[MAXN], sub[MAXN];
void pull(int x){
    sum[x] = val[x] + sum[ch[x][0]] + sum[ch[x][1]];        // PATH fold (no vir)
    sub[x] = val[x] + vir[x] + sub[ch[x][0]] + sub[ch[x][1]];// SUBTREE fold (+ vir)
}
void access(int x){
    for(int last=0; x; last=x, x=fa[x]){
        splay(x);
        int r = ch[x][1];                    // old solid (deeper) child -> becomes virtual
        if(r) vir[x] += sub[r];
        if(last) vir[x] -= sub[last];        // new solid child was virtual until now
        ch[x][1] = last;
        pull(x);
    }
}
long long subtree_sum(int x){ access(x); return sub[x]; }   // x's aux root holds all
\end{lstlisting}

\begin{insightbox}[Why the path fold ignores \(\mathit{vir}\) but the subtree fold needs it]
A path \(u\leadsto v\) is a single preferred chain -- the vertices on it are exactly
the splay subtree, and their off-path branches are NOT part of the path, so
\(\mathit{sum}\) must exclude \(\mathit{vir}\). A subtree, by contrast, is ``this
node and everything below it,'' which is precisely the branches \(\mathit{vir}\)
holds. Keeping the two folds as sibling fields in \code{pull} -- never mixing them --
is why the same structure can answer a path query and a subtree query without a
rebuild; the Summit will use only the path fold, but knowing the sibling makes that
choice deliberate rather than lucky.
\end{insightbox}

\subsection*{The connectivity question, and the LCT-vs-Euler-tour pick}

Not every dynamic-forest problem needs a path aggregate. If the only query is
``are \(u,v\) connected?'' (or ``is this graph still a forest after adding edge
\((u,v)\)?''), an LCT is overkill -- you pay splay constants for a boolean. An
Euler-tour tree stores the tour in a balanced BST and splits/merges on
\code{link}/\code{cut}, answering connectivity by testing whether two nodes lie in
one tour; a rollback DSU does it offline in near-linear time with a fraction of the
code. The decision is the pick table below.

\begin{definition}{Which dynamic-forest tool? (pick table)}{c23-map}
\emph{Path value needed online} (sum/min/max over the route) \(\Rightarrow\) full LCT
with \(\mathit{sum}\) (\(\secref{c23-makeroot}\)). \emph{Subtree value needed online}
\(\Rightarrow\) LCT with the \(\mathit{vir}\) sibling (\(\thmref{c23-vsum}\)). \emph{Only
connectivity / component identity, online} \(\Rightarrow\) Euler-tour tree (simpler
than LCT, \(\chref{c24}\)). \emph{Only connectivity, all edits known offline}
\(\Rightarrow\) rollback DSU or segment-tree-on-time (\(\chref{c28}\), engine
\(\chref{c18}\)). \emph{Fixed tree, path/subtree ranges} \(\Rightarrow\) not dynamic at
all: HLD (\(\chref{c16}\)) or Euler+Fenwick (\(\chref{c24}\)) -- see
\(\secref{c23-boundary}\).
\end{definition}

\begin{warningbox}[Subtree aggregates are where LCTs go to die]
The \(\mathit{vir}\) update in \(\coderef{c23-subtree}\) is easy to get subtly wrong:
forgetting to \emph{subtract} \(\mathit{sub}[\mathit{last}]\) (or to add
\(\mathit{sub}[r]\)) leaves a stale virtual sum that only shows up on subtree queries
whose access path demoted a child -- a rare, order-dependent WA. If your problem can
be solved by an Euler-tour tree or offline DSU, take that exit; the \(\mathit{vir}\)
discipline is the single hardest part of a production LCT and is exactly what
\(\secref{c23-boundary}\) tells you to avoid when a simpler tool fits.
\end{warningbox}

\begin{hardeningbox}
\begin{itemize}[nosep,leftmargin=1.4em]
  \item \textbf{Two folds, two uses.} Never read \(\mathit{sub}\) for a path query or
        \(\mathit{sum}\) for a subtree query; \(\mathit{sum}\) omits \(\mathit{vir}\),
        \(\mathit{sub}\) includes it. Mixing them silently mis-counts every branch.
  \item \textbf{Order inside access.} Compute \(\mathit{sub}[r]\) and
        \(\mathit{sub}[\mathit{last}]\) \emph{before} re-pointing \(\mathit{ch}[x][1]\),
        then \code{pull(x)} after -- otherwise the subtracted value is already stale.
  \item \textbf{Rooted vs rooted-after-evert.} \(\mathit{vir}\) is defined against the
        \emph{current} represented root; after \code{makeroot} the parent direction
        flips and every \(\mathit{vir}\) along the everted path changes meaning -- an
        evert plus a subtree query is legal but must recompute, not reuse cached
        \(\mathit{sub}\) from before the flip.
\end{itemize}
\end{hardeningbox}

\begin{recallbox}[25.4 in one breath]
\emph{A splay only sees solid (preferred) children; the virtual children live in
\(\mathit{vir}[x]\). \code{access} swaps a child between the two sets, so update
\(\mathit{vir}\) by moving \(\mathit{sub}[\cdot]\) values in and out
(\(\thmref{c23-vsum}\)); the subtree fold adds \(\mathit{vir}\), the path fold omits it.
And before any of this: is a path/subtree value actually required? If only
connectivity is, an Euler-tour tree or rollback DSU wins (\(\defref{c23-map}\)); if
the tree never changes, it was never an LCT problem.}
\selftest{
  \item Why does the path fold exclude \(\mathit{vir}\) but the subtree fold include it? \emph{(A path is one solid chain; a subtree includes the virtual branches, Def.~\defref{c23-virtual}.)}
  \item During \code{access}, what happens to the old right child \(r\) and the new child \(\mathit{last}\)? \emph{(\(r\) becomes virtual (+\(\mathit{sub}[r]\)), \(\mathit{last}\) becomes solid (\(-\mathit{sub}[\mathit{last}]\)), \(\thmref{c23-vsum}\).)}
  \item Name the two easier tools for connectivity-only dynamic forests. \emph{(Euler-tour tree; rollback DSU / offline D\&C, \(\defref{c23-map}\).)}
}
\end{recallbox}

% =====================================================================
%  25.5  The Engine Boundary -- when a link-cut tree is the WRONG tool.
%  Exports: sec:c23-boundary  def:c23-boundary-pick
%  WORKED examples only -- NO solvedproblem (in-section solved stays 2).
%  Discharges the forward-cites promised by Ch.22 (sec:c16-boundary) and
%  Ch.24 (link-cut is the dynamic-topology neighbour). Mirrors the
%  ch18 05_boundary worked-only style.
% =====================================================================
\section{The Engine Boundary: When Link-Cut Is the Wrong Tool}\label{sec:c23-boundary}

\begin{previouslybox}[Every engine in this book ends by naming its own limits]
Chapter~22 closes its HLD chapter by listing the problems HLD must hand away
(\(\secref{c16-boundary}\)); Chapter~24 does the same for centroid decomposition.
An LCT is the most general of the three -- it can reroot, rewire and fold a moving
path -- and generality is exactly why it is the most expensive to write and the
easiest to get wrong. The mature reflex is not ``can an LCT solve this?'' (it can
solve almost anything about a forest) but ``should I use one?'' This section is a
worked reading of six specs and the cheaper tool each one actually wants.
\end{previouslybox}

\begin{topicintro}
A link-cut tree earns its constants only when the \emph{topology} changes between
queries. Fix the tree and every job it does has a cheaper specialist; keep only
connectivity and a DSU beats it outright; know all edits offline and a rollback DSU
or a divide-and-conquer over the timeline runs in a fraction of the code. The
section walks six short specifications through the pick
(\(\defref{c23-boundary-pick}\)); none has a solved problem, because the whole
lesson is \emph{not} writing an LCT.

\textbf{Prerequisites.} The whole of \(\secref{c23-splay}\)--\(\secref{c23-beyond}\),
so you know what an LCT is before you decide against it; \(\secref{c16-boundary}\)
(HLD's boundary), Ch.~5 (\chref{c24}) Euler tour, Ch.~17 (\chref{c28}) DSU, Ch.~24
(\chref{c18}) offline D\&C.

\textbf{After this section you can}
\begin{itemize}[nosep]
  \item read a dynamic-tree spec and spot whether the topology truly changes; \hfill\lv{5}
  \item route fixed-tree path/subtree jobs to HLD / Euler+Fenwick; \hfill\lv{5}
  \item route offline connectivity to rollback DSU and reserve the LCT for online path work. \hfill\lv{5}
\end{itemize}
\end{topicintro}

\begin{definition}{The dynamic-tree tool pick}{c23-boundary-pick}
Ask two questions in order. \textbf{(Q1) Does the topology change between
queries?} No \(\Rightarrow\) it is a \emph{static} tree: path range add/sum \(\to\)
HLD (\chref{c16}); subtree range \(\to\) Euler tour + Fenwick (\chref{c24}); a value
per node \(\to\) tree DP / rerooting (\chref{c10}, \chref{c20}). \textbf{(Q2) If it
does change, do you need a value \emph{along the path}, or only
\emph{connectivity}?} Path value, online \(\Rightarrow\) LCT
(\(\secref{c23-makeroot}\)); subtree value, online \(\Rightarrow\) LCT with virtual
children (\(\secref{c23-beyond}\)) or an Euler-tour tree; connectivity only, online
\(\Rightarrow\) Euler-tour tree; connectivity only, offline \(\Rightarrow\) rollback
DSU / segment-tree-on-time (\(\chref{c28}\), \(\chref{c18}\)).
\end{definition}

\begin{comparisonbox}[Same problem, three budgets: LCT vs HLD vs offline]
\setlength{\tabcolsep}{3pt}\renewcommand{\arraystretch}{1.2}%
\begin{tabular}{@{}p{0.29\linewidth}p{0.21\linewidth}p{0.21\linewidth}p{0.21\linewidth}@{}}
\toprule
\textbf{Job} & \textbf{Fixed tree} & \textbf{Changing, online} & \textbf{Changing, offline}\\
\midrule
path sum/min & HLD + segtree \(\bigO{\log^2 n}\) & LCT \(\bigO{\log n}\) amort & offline D\&C on time \\
subtree sum & Euler + Fenwick \(\bigO{\log n}\) & ETT or LCT + \(\mathit{vir}\) & rollback / D\&C \\
connectivity & trivial (one tree) & ETT or LCT \code{findroot} & rollback DSU \(\bigO{\alpha}\) \\
\bottomrule
\end{tabular}
\end{comparisonbox}

\subsection*{Six specs, read and routed}

\begin{itemize}[leftmargin=1.4em]
  \item \textbf{``Fixed tree, path add and path sum, \(2\cdot10^5\) queries.''} No
        edge ever changes. \(\Rightarrow\) \emph{HLD} (\chref{c16}) -- an LCT would
        pass but with a larger constant and triple the code. \textbf{Do not.}
  \item \textbf{``Fixed tree, subtree update and subtree sum.''} \(\Rightarrow\)
        \emph{Euler tour + Fenwick} (\chref{c24}) in one line of insight. \textbf{Do not.}
  \item \textbf{``Add/remove edges; only ask if \(u,v\) connected; all queries known
        in advance.''} \(\Rightarrow\) \emph{rollback DSU} or D\&C on the edge-activity
        intervals (\chref{c18}). \textbf{Do not.}
  \item \textbf{``Link/cut online; only connectivity.''} \(\Rightarrow\) \emph{Euler-tour
        tree} (or an LCT \code{findroot}); the ETT is less error-prone when no path
        value is needed. \textbf{Prefer ETT.}
  \item \textbf{``Link/cut online; sum along the current path.''} \(\Rightarrow\) \textbf{LCT.}
        The genuine, textbook use (\(\secref{c23-makeroot}\)).
  \item \textbf{``Over all simple paths in a changing \emph{graph}, the bottleneck
        value.''} \(\Rightarrow\) \textbf{LCT}, but only after a condensation (block-cut
        or spanning-forest) turns it into a forest -- which is precisely the Summit.
\end{itemize}

\begin{insightbox}[The honest cost of generality]
An LCT's \code{access}/\code{splay} hides a big constant and, worse, a large
\emph{bug} surface: lazy reverse tags, virtual-child bookkeeping, \code{pushdown}
order. A fixed-tree HLD gives you the same asymptotics on \emph{no} topology change
with array-based segment trees you can debug in minutes. The reason this chapter
exists is the fifth and sixth specs -- the ones HLD genuinely cannot serve. Treat a
topology that never changes as the loudest possible signal that you wanted a static
engine.
\end{insightbox}

\begin{warningbox}[The four ways students reach for an LCT too early]
\begin{enumerate}[nosep,leftmargin=1.6em]
  \item ``the problem mentions a dynamic tree'' but the edges are fixed and only
        \emph{vertex values} change -- that is HLD/Euler, not LCT.
  \item connectivity-only, asked offline -- a rollback DSU is \(\bigO{\alpha}\) with a
        tenth of the code (\chref{c28}).
  \item counting over \emph{all} paths of a fixed tree -- centroid decomposition
        (\chref{c18}), not a per-path fold.
  \item a state per node with no route dimension -- tree DP / rerooting (\chref{c10}).
\end{enumerate}
Each wastes an LCT on a job a simpler engine does faster and correctly.
\end{warningbox}

\begin{hardeningbox}
\begin{itemize}[nosep,leftmargin=1.4em]
  \item \textbf{Prove the topology moves.} Before coding, find the exact operation that
        adds or deletes an \emph{edge}. If you cannot, you have an HLD/Fenwick problem.
  \item \textbf{Check the constant.} With \(n,q\le 2\cdot10^5\) and a fixed tree, an
        LCT's \(\sim\)20 splay-steps per query versus HLD's segment-tree walk can be the
        difference between AC and TLE even though both are ``\(O(\log n)\).''
  \item \textbf{Offline \(\Rightarrow\) no online structure.} If you may read all queries
        first, prefer rollback DSU / D\&C; the LCT is the tool for when you \emph{cannot}.
\end{itemize}
\end{hardeningbox}

\begin{recallbox}[25.5 in one breath]
\emph{An LCT is for a \textbf{changing topology} with a \textbf{path (or subtree) value}
query, \textbf{online}. Fixed tree \(\to\) HLD / Euler+Fenwick; connectivity only \(\to\)
ETT or rollback DSU; all-pairs counting \(\to\) centroid; per-node state \(\to\) tree DP.
Two questions decide it: does an edge ever appear or vanish, and do you need a value
along the route? Both ``yes'' \(\Rightarrow\) LCT; otherwise hand it off
(\(\defref{c23-boundary-pick}\)).}
\selftest{
  \item A spec has path sums but the tree never changes -- which engine? \emph{(HLD, \chref{c16}; the LCT's generality is dead weight.)}
  \item Connectivity is all you need and queries are offline -- best tool? \emph{(Rollback DSU / D\&C on time, \chref{c28}/\chref{c18}.)}
  \item State the two questions of the pick. \emph{(Does the topology change between queries? Do you need a value along the path? \(\defref{c23-boundary-pick}\).)}
}
\end{recallbox}

% =====================================================================
%  Chapter 25 review -- the 14-station revision contract.
%  Revision-only: tables, one-liners, pointers back via \secref/\thmref/
%  \defref/\figref/\coderef/\probref.  No re-derivations.
%  Summit: CF 487E "Tourists" (3200, L6) -- a connected undirected graph with a
%  price w_i per city and two online queries: C a w changes city a's price, and
%  A a b asks for the cheapest souvenir obtainable on SOME simple a-b route.
%  A simple route can detour inside any biconnected block it enters, so the
%  answer is a PATH-MINIMUM on the block-cut tree where each block node stores
%  the min price of its vertices and each vertex node stores its own price.
%  Condense once (Tarjan BCT), then run the chapter's LCT path-min; a price
%  change re-pushes the multiset minima of the (few) blocks containing the city.
%  Both official samples hand-replayed: S1 -> 1, 2 ; S2 -> 2, 1, 5, 3.
%  Exports prob:c23-summit, lst:c23-kit, lst:c23-summit.
%  Mixed set is EXACTLY 10, all live-verified (names/ratings/tags re-pulled from
%  the CF API this session), un-owned, sorted ASCENDING by rating; each hint
%  names the reduction (or the REJECTION) and cites a c23 listing.  CF 1601C is
%  OWNED by Ch.19 -- never linked here.
% =====================================================================

\chaptersection{Chapter 25 Summary}

One sentence: \emph{a splay tree is a binary search tree that shuffles whatever you
just touched to the root, so its \emph{amortised} cost tracks recent use rather than
worst-case depth; a link-cut tree strings those splays along \emph{preferred paths}
and re-preferences them with a single \code{access} loop, which is what lets one
structure reroot, rewire, and fold an aggregate over the current path of a forest
whose edges keep changing (\(\secref{c23-splay}\), \(\secref{c23-makeroot}\)).}
Chapter~22 could only fold a path in a \emph{fixed} tree; this chapter drops that.
Each row below is a pointer, not a proof.

\begin{center}
\renewcommand{\arraystretch}{1.25}
\begin{tabular}{@{}p{0.24\linewidth} p{0.48\linewidth} p{0.20\linewidth}@{}}
\toprule
\textbf{Result} & \textbf{Statement} & \textbf{Where} \\
\midrule
Splay & Repeatedly rotate the accessed key to the root (zig, zig-zig, zig-zag); with
potential \(\Phi=\sum_x\log\size(x)\) every operation is \(\bigO{\log n}\) amortised. &
\defref{c23-splay}, \thmref{c23-splay-amort} \\
Preferred path & The forest is covered by vertex-disjoint \emph{solid} paths, each held in
one aux splay; a non-aux child carries a \emph{path-parent} pointer, so splay-parent and
path-parent are two different roles of \code{fa}. & \defref{c23-prefer}, \figref{c23-aux} \\
Access & \code{access(x)} walks \(x\) to the represented root, at each step splaying and
re-pointing the right child to the last node visited; this makes root\(\to x\) one solid
path and touches \(\bigO{\log n}\) splays amortised. & \defref{c23-access}, \coderef{c23-lct-core} \\
Two-access LCA & \code{access(u)} then climbing \(v\) through path-parents with a splay at
each hop records the last aux root touched -- that node is \(\lca(u,v)\) in a rooted forest. &
\thmref{c23-access-lca}, \coderef{c23-dynalca} \\
Evert / path fold & \code{makeroot(x)} is \code{access(x)} plus a reverse lazy tag on the
whole aux tree; \code{makeroot(u); access(v)} exposes the \(u\)--\(v\) path as \(v\)'s splay,
so any fold monoid reads off \(\mathit{sum}[v]\). & \defref{c23-makeroot}, \thmref{c23-path-query} \\
Subtree fold & A splay sees only \emph{solid} children; the virtual children live in a
separate accumulator \(\mathit{vir}[x]\) that \code{access} updates as it swaps a child
between solid and virtual -- the path fold omits it, the subtree fold includes it. &
\defref{c23-virtual}, \thmref{c23-vsum} \\
Engine map & A link-cut tree earns its constants only when the \emph{topology} moves and a
\emph{path (or subtree) value} is needed online; connectivity goes to ETT/DSU, fixed paths to
HLD, all-pairs to centroid. & \defref{c23-map}, \defref{c23-boundary-pick} \\
\bottomrule
\end{tabular}
\end{center}

\medskip
{\small\textbf{Algorithms from this chapter.}}
\begin{center}
\renewcommand{\arraystretch}{1.2}
\begin{tabular}{@{}p{0.36\linewidth} l l >{\raggedright\arraybackslash}p{0.20\linewidth}@{}}
\toprule
\textbf{Algorithm} & \textbf{Prep} & \textbf{Per op} & \textbf{Code} \\
\midrule
Bottom-up splay (rotate to root) & -- & \(\bigO{\log n}\) amort & \coderef{c23-splay-rot} \\
Top-down splay (no recursion/parents) & -- & \(\bigO{\log n}\) amort & \coderef{c23-topdown} \\
SPOJ ORDERSET (order-statistics set) & \(\bigO{1}\) & \(\bigO{\log n}\) amort & \coderef{c23-orderset} \\
LCT core: access / findroot (rooted) & \(\bigO{n}\) & \(\bigO{\log n}\) amort & \coderef{c23-lct-core} \\
SPOJ DYNALCA (link/cut/lca) & \(\bigO{n}\) & \(\bigO{\log n}\) amort & \coderef{c23-dynalca} \\
Evert + path-sum LCT & \(\bigO{n}\) & \(\bigO{\log n}\) amort & \coderef{c23-evert}, \coderef{c23-pathsum} \\
Subtree-sum LCT (\(\mathit{vir}\) sibling) & \(\bigO{n}\) & \(\bigO{\log n}\) amort & \coderef{c23-subtree} \\
Engine pick table (is it even LCT?) & -- & -- & \defref{c23-boundary-pick}, \defref{c23-map} \\
Link-cut kit (one listing) & -- & -- & \coderef{c23-kit} \\
CF 487E (Summit) -- BCT + LCT path-min & \(\bigO{n+m}\) & \(\bigO{\log n}\) amort & \coderef{c23-summit} \\
\bottomrule
\end{tabular}
\end{center}

\chaptersection{Master Decision Flowchart}

Post-study ladder: walk top to bottom, stop at the first branch that matches. It decides
\emph{whether the topology truly moves} (the LCT's only excuse for existing) and, if so,
\emph{which operation} -- \code{access} alone, \code{makeroot} plus a path fold, or the
\(\mathit{vir}\) subtree sibling -- the query needs. Everything that fails the top gate is
handed back to a static engine (HLD \(\chref{c16}\), Euler+Fenwick \(\chref{c24}\), offline DSU
\(\chref{c28}\)). A \emph{graph} that is not a forest must be \emph{condensed} first -- the Summit.

\begin{lstlisting}[language={}, basicstyle=\ttfamily\small, numbers=none, frame=single]
  Does the structure answer queries about a FOREST WHOSE EDGES CHANGE (link / cut)?
  |
  |-- The tree is FIXED -- no edge is ever added or deleted, only VALUES move?
  |     +-- YES: NOT link-cut (def:c23-boundary-pick).  Path range add/sum -> HLD
  |         (Ch.22, c16); subtree range -> Euler tour + Fenwick (Ch.5, c24).
  |         Reaching for an LCT here pays splay constants for nothing.
  |
  |-- Only CONNECTIVITY (are u,v in one tree?) -- never a path value?
  |     +-- offline (all edits known first): rollback DSU / segment-tree-on-time
  |         (Ch.17 c28, Ch.24 c18).  Fastest, least code.  NOT LCT.
  |     +-- online: Euler-tour tree (Ch.5, c24) or an LCT findroot -- but if no
  |         path value is needed, the ETT is the smaller, safer structure.
  |
  |-- A PATH (or SUBTREE) VALUE on a forest whose edges DO change, online: THIS chapter.
  |     First own the splay engine (25.1): rot / splay, zig-zig vs zig-zag, size
  |     fields for k-th.  Without the amortised argument (thm:c23-splay-amort) the
  |     LCT is just a slow BST.  lst:c23-splay-rot, lst:c23-topdown.
  |
  |     |-- ROOTED growing forest, need LCA / root / connected?  (25.2)
  |     |     access(x) exposes root->x; findroot = access then leftmost then splay;
  |     |     lca(u,v) = access(u) then the two-access climb on v.  lst:c23-lct-core,
  |     |     lst:c23-dynalca.  NO reverse tag here -- keep the root fixed.
  |     |
  |     |-- ARBITRARY path u~v fold (sum/min/max) with rerooting?  (25.3)
  |     |     makeroot(x)=access+apply-reverse; query = makeroot(u); access(v);
  |     |     read sum[v].  link = makeroot(u); fa[u]=v.  cut = makeroot(u);
  |     |     access(v); detach.  Push rev tags top-down BEFORE any rotate.
  |     |     lst:c23-evert, lst:c23-pathsum.
  |     |
  |     |-- SUBTREE aggregate on a moving tree?  (25.4)
  |     |     add the vir[] sibling: access moves a child between solid and virtual,
  |     |     so vir[x]-=sub[child-in], vir[x]+=sub[old-right], pull after.  sub[]
  |     |     folds val+vir+solid; the PATH fold sum[] must NOT include vir.
  |     |     lst:c23-subtree.  (thm:c23-vsum)
  |     |
  |     |-- The input is a GRAPH, not a forest?  Condense, THEN run the LCT.
  |         Block-cut (biconnected) tree or a maintained spanning forest turns the
  |         query into a path fold on a tree -- this is exactly the Summit, CF 487E:
  |         "cheapest city on SOME simple route" -> path-min on the block-cut tree
  |         where a block node carries the min price of its vertices.  lst:c23-summit.
  |
  |-- NONE of the above -- you wanted a static engine.  See 25.5 (def:c23-boundary-pick):
        fixed tree -> HLD / Euler; counting over ALL paths -> centroid (Ch.24, c18);
        per-node state -> tree DP (Ch.13, c10); offline node-sets -> Mo (Ch.27, c21).
\end{lstlisting}

\chaptersection{Common Mistakes Checklist}

\begin{itemize}[nosep, leftmargin=1.4em]
  \item \textbf{Confusing the splay-parent with the path-parent.} \code{fa[x]} is a splay edge
        when \(x\) is an aux child and a \emph{path-parent} when \(x\) is an aux root; the test
        is \(\mathit{isRoot}(x)\): \code{fa}==0 \emph{or} \code{fa} has no child equal to \(x\).
        Writing \(\mathit{isRoot}(x)\equiv(\mathit{fa}[x]=0)\) breaks every access. \(\defref{c23-prefer}\), \(\coderef{c23-lct-core}\).
  \item \textbf{Not pushing the reverse tag top-down before rotating.} After \code{makeroot} the
        whole aux subtree is lazily flipped; a \code{splay} that rotates before pushing \code{rev}
        to the ancestors reattaches mirrored children -- \textbf{WA} that only shows on deep
        paths. Push from the aux root down before the zig-zig loop. \(\defref{c23-makeroot}\), \(\coderef{c23-evert}\).
  \item \textbf{Forgetting \code{pull} after \code{access} re-points the right child.} The line
        \(\mathit{ch}[x][1]\leftarrow\mathit{last}\) changes the splay subtree, so \(\mathit{sum}\)/\(\mathit{mn}\)
        must be recomputed before the next hop, or the aggregate is stale. \(\defref{c23-access}\), \(\coderef{c23-lct-core}\).
  \item \textbf{\code{link} without \code{makeroot}.} \code{link(u,v)} must expose \(u\) as a
        represented root first (\code{makeroot(u); fa[u]=v}); linking a node that already has a
        parent silently welds two trees into a structure that is not a forest. \(\coderef{c23-evert}\).
  \item \textbf{\code{findroot} not splaying at the end.} Walking to the leftmost node then
        \emph{not} splaying it destroys the amortised bound and leaves the exposed root buried;
        always \code{splay} the returned root. \(\coderef{c23-lct-core}\).
  \item \textbf{Reading the subtree fold off the path aggregate (or vice-versa).} \(\mathit{sum}\)
        omits \(\mathit{vir}\); \(\mathit{sub}\) includes it. A subtree query answered from
        \(\mathit{sum}\), or a path query from \(\mathit{sub}\), mis-counts every branch. \(\thmref{c23-vsum}\), \(\coderef{c23-subtree}\).
  \item \textbf{Wrong \(\mathit{vir}\) update order in \code{access}.} Compute
        \(\mathit{sub}[\mathit{last}]\) and \(\mathit{sub}[r]\) \emph{before} re-pointing
        \(\mathit{ch}[x][1]\); after the re-point the demoted child's value is unreachable, so a
        late read returns garbage. \(\thmref{c23-vsum}\), \(\coderef{c23-subtree}\).
  \item \textbf{Using an LCT on a tree that never changes.} If no edge is ever added or deleted,
        HLD has the same asymptotics with a debuggable array segment tree; the LCT is pure
        overhead and a larger bug surface. \(\defref{c23-boundary-pick}\), \(\secref{c23-boundary}\).
  \item \textbf{Sentinel node \(0\) not neutral.} Node \(0\) is the null: its \(\mathit{size}=0\)
        for order statistics and its \(\mathit{mn}=+\infty\) for a path minimum; forget the
        neutral value and the first rotation corrupts the aggregate. \(\coderef{c23-orderset}\), \(\coderef{c23-summit}\).
\end{itemize}

\chaptersection{WA Clinic}

Verdicts are symptoms; the causes are the short list this chapter drilled. Diagnose top-down --
every leaf names a listing.

\begin{figure}[H]
\centering
\begin{tikzpicture}[diagram,
    v/.style={boxnode, fill=primaryblue!12, draw=primaryblue, text width=22mm, align=center},
    b/.style={boxnode, fill=difforange!12, draw=difforange, text width=12mm},
    l/.style={boxnode, fill=successgreen!10, draw=successgreen, text width=48mm,
              font=\scriptsize, align=left}]
  \node[v] (verdict) {\textbf{Verdict on the judge}};
  \node[b, below left=8mm and 5mm of verdict] (wa) {\textbf{WA}};
  \node[b, below right=8mm and 5mm of verdict] (tle) {\textbf{TLE / RE}};
  \node[l, below right=1mm and 2mm of wa.south] (wa1)
    {the reverse tag was never pushed before a rotate -- deep paths come back mirrored, so a
     path fold reads the wrong vertices (\(\defref{c23-makeroot}\), \(\coderef{c23-evert}\))};
  \node[l, below left=1mm and 2mm of wa.south] (wa2)
    {\(\mathit{isRoot}\) written as \code{fa==0}, conflating path-parent with splay-parent, or a
     subtree query answered from the path fold (\(\coderef{c23-lct-core}\), \(\thmref{c23-vsum}\))};
  \node[l, below right=1mm and 2mm of tle.south] (tle1)
    {no amortised discipline: \code{findroot} fails to splay, or the tree never changes so an LCT
     is paying splay constants where HLD was free (\(\coderef{c23-lct-core}\), \(\defref{c23-boundary-pick}\))};
  \node[l, below left=1mm and 2mm of tle.south] (tle2)
    {recursion depth on a chain splay blows the stack, or the sentinel \(0\) carries a non-neutral
     aggregate (\(\coderef{c23-topdown}\), \(\coderef{c23-summit}\))};
  \draw[flow] (verdict) -- (wa);
  \draw[flow] (verdict) -- (tle);
  \draw[flow] (wa) -- (wa1);
  \draw[flow] (wa) -- (wa2);
  \draw[flow] (tle) -- (tle1);
  \draw[flow] (tle) -- (tle2);
\end{tikzpicture}
\caption*{WA: the lazy reverse tag is not pushed before rotating, or \(\mathit{isRoot}\) mistakes a
path-parent for a splay-parent, or the path/subtree folds are swapped. TLE or RE: \code{findroot}
forgets its final splay, the topology never actually changes (use HLD), or the null sentinel is not
neutral so the first rotate corrupts an aggregate.}
\end{figure}

\chaptersection{Quick Reference Card}

\begin{quickref}[Chapter 25 at a glance]
\textbf{A link-cut tree is a set of splay trees, one per preferred path of a forest whose edges
change; \code{access} re-preferences the root-to-node path in one loop, \code{makeroot} reroots with
a lazy flip, and any path (or, with \(\mathit{vir}\), subtree) fold answers in \(\bigO{\log n}\) amortised. }

\medskip
\textbf{Splay (25.1).} Rotate the touched key to the aux root: zig (root child), zig-zig (two same-side,
one rotate of the parent first), zig-zag (opposite sides, two rotates of \(x\)). With
\(\Phi=\sum_x\log\size(x)\) this is \(\bigO{\log n}\) amortised; keep \(\mathit{size}\) for k-th.
(\(\defref{c23-splay}\), \(\thmref{c23-splay-amort}\), \(\coderef{c23-splay-rot}\), \(\coderef{c23-topdown}\)).

\medskip
\textbf{Access / rooted (25.2).} \(\mathit{isRoot}(x)\) is a splay test, not \code{fa==0}. \code{access(x)}:
splay, set right child to \code{last}, pull, step to \code{fa}. \code{findroot}: access then leftmost then
splay. \code{lca(u,v)}: access(u); climb \(v\) by path-parent splaying, return last.
(\(\defref{c23-access}\), \(\thmref{c23-access-lca}\), \(\coderef{c23-lct-core}\), \(\coderef{c23-dynalca}\)).

\medskip
\textbf{Evert / paths (25.3).} \code{makeroot(x)}=access+flip (swap children, toggle rev). \code{query(u,v)}=
makeroot(u); access(v); read \(\mathit{fold}[v]\). \code{link}: makeroot(u); fa[u]=v. \code{cut}: makeroot(u);
access(v); detach \(v\)'s left child. Push rev top-down before rotating.
(\(\defref{c23-makeroot}\), \(\thmref{c23-path-query}\), \(\coderef{c23-evert}\), \(\coderef{c23-pathsum}\)).

\medskip
\textbf{Subtree / boundary (25.4).} Keep \(\mathit{vir}[x]\) = fold over virtual children; \code{access}
moves a child between the solid splay and \(\mathit{vir}\) (subtract the promoted, add the demoted).
\(\mathit{sub}=\mathit{val}+\mathit{vir}+\) solid, \(\mathit{sum}\) omits \(\mathit{vir}\).
(\(\defref{c23-virtual}\), \(\thmref{c23-vsum}\), \(\coderef{c23-subtree}\)).

\medskip
\textbf{The boundary (25.5).} Fixed tree \(\to\) HLD (\(\chref{c16}\)) / Euler+Fenwick (\(\chref{c24}\));
connectivity only, offline \(\to\) rollback DSU (\(\chref{c28}\)); connectivity only, online \(\to\)
Euler-tour tree; all-pairs counting \(\to\) centroid (\(\chref{c18}\)); \emph{moving forest + online path
value} \(\to\) LCT (\(\defref{c23-boundary-pick}\), \(\defref{c23-map}\)).
\end{quickref}

\begin{lstlisting}[caption={Chapter 25 kit: the splay rotate-and-splay, the rooted LCT core (isRoot, access, findroot), evert with its reverse tag plus link/cut/path-query, the subtree fold with its vir sibling, the two-access LCA, and a path-min variant -- one compact skeleton each, in one place. ASCII only, identifiers spelled without underscores.}, label={lst:c23-kit}]
// ===== SPLAY (25.1) =====
int ch[MAXN][2], fa[MAXN];
bool isRoot(int x){ int f=fa[x]; return f==0 || (ch[f][0]!=x && ch[f][1]!=x); }
void rot(int x){ int y=fa[x], z=fa[y], d=(ch[y][1]==x);
    if(!isRoot(y)) ch[z][ch[z][1]==y]=x; fa[x]=z;
    ch[y][d]=ch[x][d^1]; if(ch[x][d^1]) fa[ch[x][d^1]]=y;
    ch[x][d^1]=y; fa[y]=x; pull(y); pull(x); }
void splay(int x){ pushPath(x);
    while(!isRoot(x)){ int y=fa[x], z=fa[y];
        if(!isRoot(y)) rot( (ch[y][0]==x)==(ch[z][0]==y) ? y : x );  // zig-zig vs zig-zag
        rot(x); } }

// ===== ROOTED LCT CORE (25.2) : NO reverse tag, root is fixed =====
void access(int x){ for(int last=0; x; last=x, x=fa[x]){ splay(x); ch[x][1]=last; pull(x);} }
int findroot(int x){ access(x); while(ch[x][0]) { push(x); x=ch[x][0]; } splay(x); return x; }

// ===== EVERT + PATH FOLD (25.3) : sum is a generic monoid fold =====
bool rev[MAXN];
void pushRev(int x){ if(!x) return; swap(ch[x][0],ch[x][1]); rev[x]^=1; }
void push(int x){ if(rev[x]){ pushRev(ch[x][0]); pushRev(ch[x][1]); rev[x]=0; } }
void makeroot(int x){ access(x); pushRev(x); }
void link(int u,int v){ makeroot(u); fa[u]=v; }
void cut(int u,int v){ makeroot(u); access(v); if(ch[v][0]==u && !ch[u][1]){ ch[v][0]=0; fa[u]=0; pull(v);} }
ll query(int u,int v){ makeroot(u); access(v); return sum[v]; }   // fold over u~v

// ===== SUBTREE FOLD with vir sibling (25.4) =====
void accessV(int x){ for(int last=0; x; last=x, x=fa[x]){ splay(x);
    int r=ch[x][1];                                        // old solid child -> virtual
    if(r) vir[x]+=sub[r]; if(last) vir[x]-=sub[last];
    ch[x][1]=last; pull(x); } }                            // pull: sum=val+solid ; sub=val+vir+solid
ll subtree(int x){ accessV(x); return sub[x]; }

// ===== TWO-ACCESS LCA (25.2, rooted) =====
int lca(int u,int v){ if(findroot(u)!=findroot(v)) return 0;
    access(u); int last=0; for(int x=v; x; x=fa[x]){ splay(x); last=x; } splay(last); return last; }

// ===== PATH-MIN variant (Summit 487E) : mn[0]=+INF, mn=min(val, children) =====
void pullMin(int x){ mn[x]=val[x]; if(ch[x][0]) mn[x]=std::min(mn[x],mn[ch[x][0]]);
    if(ch[x][1]) mn[x]=std::min(mn[x],mn[ch[x][1]]); }
ll pathMin(int u,int v){ makeroot(u); access(v); return mn[v]; }
\end{lstlisting}

\chaptersection{Write-From-Memory Drills}

Five signature skeletons, each lifted from a listing already in the chapter. Cover the marked line,
reproduce from memory, then check the answer one line below.

\begin{lstlisting}[caption={Drill 1 -- splay step: decide zig-zig vs zig-zag, then rotate (from lst:c23-splay-rot)}]
void splay(int x){ pushPath(x);
    while(!isRoot(x)){
        int y = fa[x], z = fa[y];
        if(!isRoot(y))                       // grandparent exists -> two-level case
            rot( (ch[y][0]==x) == (ch[z][0]==y) ? y : x );   // same side: rotate y; else rotate x
        rot(x);                              // always finish by rotating x
    }
}
\end{lstlisting}
\noindent{\scriptsize Key: the parity test \((\mathit{ch}[y][0]{=}x)==(\mathit{ch}[z][0]{=}y)\) is
zig-zig (grandparent and parent go the same way, rotate the \emph{parent} first) versus zig-zag
(opposite ways, rotate \emph{x} first); push the lazy tags from the aux root down before the loop
(\(\defref{c23-splay}\), \(\thmref{c23-splay-amort}\), \(\coderef{c23-splay-rot}\)).}

\begin{lstlisting}[caption={Drill 2 -- the access loop: the whole engine in four lines (from lst:c23-lct-core)}]
void access(int x){
    for(int last = 0; x; last = x, x = fa[x]){
        splay(x);            // x to the top of its aux tree
        ch[x][1] = last;     // re-prefer: the path below x becomes the right spine
        pull(x);             // refresh x's aggregate after the child swap
    }
}
\end{lstlisting}
\noindent{\scriptsize Key: after the step \(x\leftarrow\mathit{fa}[x]\), the old right child is
\emph{demoted} (it becomes virtual) and \code{last} is \emph{promoted} to solid; that single
re-point, \(\bigO{\log n}\) times, is what exposes root\(\to\)the original \(x\) as one solid path
(\(\defref{c23-access}\), \(\coderef{c23-lct-core}\)).}

\begin{lstlisting}[caption={Drill 3 -- evert and read a path fold (from lst:c23-evert)}]
void makeroot(int x){ access(x); pushRev(x); }        // x becomes the represented root
ll  query(int u,int v){ makeroot(u); access(v); return sum[v]; }  // u~v is v's aux subtree
void link(int u,int v){ makeroot(u); fa[u]=v; }       // hang u (now a root) under v
void cut(int u,int v){ makeroot(u); access(v);
    if(ch[v][0]==u && !ch[u][1]){ ch[v][0]=0; fa[u]=0; pull(v); } }
\end{lstlisting}
\noindent{\scriptsize Key: \code{makeroot} is exactly \code{access} plus a lazy reversal of the whole
aux tree; once \(u\) is the root, \code{access(v)} makes the \(u\)--\(v\) route the right spine, so
the fold sits at \(v\). The \code{cut} guard (u is v's left child with no right child) prevents
severing a non-edge (\(\defref{c23-makeroot}\), \(\thmref{c23-path-query}\), \(\coderef{c23-evert}\)).}

\begin{lstlisting}[caption={Drill 4 -- the two-access LCA on a rooted forest (from lst:c23-dynalca)}]
int lca(int u, int v){
    if(findroot(u) != findroot(v)) return 0;          // different trees -> no LCA
    access(u);                                        // expose root->u as one solid path
    int last = 0;
    for(int x = v; x; x = fa[x]) { splay(x); last = x; }  // climb v's path-parents
    splay(last);                                      // keep the amortised bound
    return last;                                      // first node already in u's path = LCA
}
\end{lstlisting}
\noindent{\scriptsize Key: after \code{access(u)} the nodes root\(\to u\) sit on solid paths; climbing
\(v\) by \code{fa} and splaying each hop, the \emph{last} aux root reached before you fall off to the
shared part is \(\lca(u,v)\) -- the two-access method, no reverse tag, root fixed
(\(\thmref{c23-access-lca}\), \(\coderef{c23-dynalca}\)).}

\begin{lstlisting}[caption={Drill 5 -- subtree fold via the vir sibling in access (from lst:c23-subtree)}]
void accessSub(int x){
    for(int last = 0; x; last = x, x = fa[x]){
        splay(x);
        int r = ch[x][1];                 // solid child about to be demoted to virtual
        if(r)    vir[x] += sub[r];        // add its WHOLE subtree into the virtual pool
        if(last) vir[x] -= sub[last];     // remove the promoted child from the pool
        ch[x][1] = last;
        pull(x);                          // sum=val+solid (no vir) ; sub=val+vir+solid
    }
}
\end{lstlisting}
\noindent{\scriptsize Key: the path fold \(\mathit{sum}\) ignores \(\mathit{vir}\) because off-path
branches are not on a path; the subtree fold \(\mathit{sub}\) must add \(\mathit{vir}\) because those
branches \emph{are} in the subtree. Read \(\mathit{sub}[\mathit{last}]\) and \(\mathit{sub}[r]\) before
re-pointing \(\mathit{ch}[x][1]\) (\(\thmref{c23-vsum}\), \(\coderef{c23-subtree}\)).}


\chaptersection{Complete Problem Checklist}

Every problem this chapter offers -- two fully solved in-section, worked traces on the running
forest \(R\), and one Summit. Tick \emph{only} after an unaided solve; a hinted attempt counts as half.

\begin{center}
\footnotesize
\setlength{\tabcolsep}{4pt}
\renewcommand{\arraystretch}{1.25}
\begin{tabular}{@{}c p{0.16\linewidth} l p{0.20\linewidth} >{\raggedright\arraybackslash}p{0.26\linewidth} c@{}}
\toprule
\textbf{\#} & \textbf{Problem} & \textbf{Source} & \textbf{Key technique} & \textbf{What to look for} & \textbf{Done?} \\
\midrule
\multicolumn{6}{l}{\textbf{Section 25.1 -- the splay engine}} \\
W25.1 & replay splay on keys \(\{1,2,3,4\}\) & worked & \lv{4} zig / zig-zig / zig-zag & access order reshapes depth; \(\mathit{size}\) kept on every rotate; \(\Phi\) drops & \(\square\) \\
W25.1b & Order Statistic Set & SPOJ & \lv{4} splay + \(\mathit{size}\) kth & dedup insert, kth by size, count below \(x\), erase; sample \(\to\) 1,2,2,invalid & \(\square\) \\
\midrule
\multicolumn{6}{l}{\textbf{Section 25.2 -- preferred paths and access}} \\
W25.2 & \code{access(7)} on forest \(R\) & worked & \lv{4} re-preference loop & root\(\to\)7 solid; siblings demoted to virtual; one aux splay per path & \(\square\) \\
W25.2b & Dynamic LCA & SPOJ & \lv{5} two-access LCA & link/cut/\emph{lca}, rooted, no reverse tag; sample \(\to\) 1,2,3,2 & \(\square\) \\
\midrule
\multicolumn{6}{l}{\textbf{Section 25.3 -- evert, link, cut, path fold (worked only)}} \\
W25.3 & path sum on chain 1--2--3--4 & worked & \lv{5} makeroot + fold & query(1,4)=10; add node \(\to\) 15; cut splits \(\{1,2\},\{3,4\}\); push rev before rotate & \(\square\) \\
\midrule
\multicolumn{6}{l}{\textbf{Section 25.4 -- virtual children and connectivity (worked only)}} \\
W25.4 & subtree sum after \code{access} & worked & \lv{5} \(\mathit{vir}\) sibling & demote solid \(\Rightarrow\) add to \(\mathit{vir}\); promote \(\Rightarrow\) subtract; \(\mathit{sub}\ne\mathit{sum}\) & \(\square\) \\
\midrule
\multicolumn{6}{l}{\textbf{Section 25.5 -- the engine boundary (worked only)}} \\
W25.5 & route six specs off the LCT & worked & \lv{5} which engine & fixed \(\to\) HLD/Euler; offline conn. \(\to\) DSU; all-pairs \(\to\) centroid; moving path \(\to\) LCT & \(\square\) \\
\midrule
\multicolumn{6}{l}{\textbf{Summit (chapter end)}} \\
W25 & Tourists & CF 487E & \lv{6} block-cut tree + LCT path-min & condense graph to BCT, block node \(=\) min price; reruns \(\to\) S1: 1,2; S2: 2,1,5,3 & \(\square\) \\
\bottomrule
\end{tabular}
\end{center}

\chaptersection{Recommended Study Order}

Five focused days. Day~1 internalises the splay rotations and ORDERSET; day~2 the access loop and
DYNALCA; day~3 evert and the path fold; day~4 the \(\mathit{vir}\) subtree sibling and the engine
boundary; day~5 the Summit and the mixed set.

\begin{center}
\small
\renewcommand{\arraystretch}{1.25}
\begin{tabular}{@{}llp{0.55\linewidth}@{}}
\toprule
\textbf{Day} & \textbf{Focus} & \textbf{Tasks} \\
\midrule
1 & Read + do 25.1 & hand-splay keys \(\{1,2,3,4\}\) after accessing 1 then 4, name each zig/zig-zig/zig-zag (\(\figref{c23-splaycases}\)); write \(\coderef{c23-splay-rot}\) cold; then SPOJ ORDERSET with the \(\mathit{size}\) kth (\(\coderef{c23-orderset}\)) \\
2 & Read + do 25.2 & trace \code{access(7)} on \(R\), marking which children become virtual (\(\coderef{c23-lct-core}\)); then SPOJ DYNALCA, replaying the sample to \(1,2,3,2\) (\(\coderef{c23-dynalca}\)) \\
3 & Read + do 25.3 & trace \code{makeroot}/\code{query} on the chain 1--2--3--4 to \(10\) then \(15\) (\(\coderef{c23-evert}\), \(\coderef{c23-pathsum}\)); write the reverse-tag pushdown cold \\
4 & Read + do 25.4 + 25.5 & add the \(\mathit{vir}\) sibling and re-trace the subtree query (\(\coderef{c23-subtree}\)); place the six specs on the pick (\(\defref{c23-boundary-pick}\)) \\
5 & Summit + mixed & condense the CF 487E sample graph to its block-cut tree and replay the answers \(2,1,5,3\) (\(\probref{c23-summit}\)); then the first 3 mixed entries \\
\bottomrule
\end{tabular}
\end{center}

\chaptersection{Difficulty Progression}

Tiers are competency gates -- each row names the \emph{skill} the tier unlocks.

\begin{center}
\small
\renewcommand{\arraystretch}{1.3}
\begin{tabular}{@{}lp{0.32\linewidth}p{0.40\linewidth}@{}}
\toprule
\textbf{Tier} & \textbf{Problems} & \textbf{What passing this tier proves} \\
\midrule
\lv{3} \ \(\bigstar\bigstar\bigstar\)
  & a single splay rotate; one \code{access} step
  & you keep the splay/path-parent distinction and refresh the aggregate after each child swap \\
\lv{4} \ \(\bigstar\bigstar\bigstar\bigstar\)
  & ORDERSET; \code{access}/\code{findroot} on \(R\)
  & you maintain \(\mathit{size}\) for k-th and expose a root-to-node path correctly \\
\lv{5} \ \(\bigstar\bigstar\bigstar\bigstar\bigstar\)
  & DYNALCA, evert path fold, \(\mathit{vir}\) subtree, mixed at CF~2100--2400
  & you push reverse tags before rotating and keep \(\mathit{sub}\) distinct from \(\mathit{sum}\) \\
\lv{6} \ \(\bigstar\bigstar\bigstar\bigstar\bigstar\bigstar\)
  & CF 487E (Summit), the CF~3000--3400 mixed set
  & you \emph{condense a graph to a block-cut tree}, then run an LCT path fold on it \\
\bottomrule
\end{tabular}
\end{center}

\chaptersection{Reduction Gallery}

Every entry is ``one dynamic-forest shape \(\to\) one LCT move (or an honest hand-off).'' Each line
collapses a problem family onto a listing already in the chapter.

\begin{itemize}[nosep, leftmargin=1.4em]
  \item ``maintain an ordered set with k-th / rank / predecessor under insert--delete'' (SPOJ ORDERSET)
        \(\to\) a splay keyed by value with a \(\mathit{size}\) field; k-th walks by subtree sizes
        (\(\defref{c23-splay}\), \(\coderef{c23-orderset}\)).
  \item ``link/cut a \emph{rooted} forest, ask LCA / connectivity online'' (SPOJ DYNALCA) \(\to\) the
        rooted core: \code{access} + the two-access climb, \code{findroot} for the same-tree test
        (\(\thmref{c23-access-lca}\), \(\coderef{c23-dynalca}\)).
  \item ``link/cut online, fold sum/min/max over the current path \(u\)--\(v\)'' \(\to\)
        \code{makeroot(u); access(v); read fold[v]}, with a lazy reverse tag (\(\thmref{c23-path-query}\),
        \(\coderef{c23-evert}\)).
  \item ``subtree aggregate on a moving tree'' \(\to\) the \(\mathit{vir}\) sibling kept current in
        \code{access}; subtree fold \(\mathit{sub}\) includes it, path fold \(\mathit{sum}\) does not
        (\(\thmref{c23-vsum}\), \(\coderef{c23-subtree}\)).
  \item ``cheapest vertex on \emph{some} simple route \(a\to b\) in a graph, with price changes''
        (CF 487E, Summit) \(\to\) condense to the block-cut tree, a block node carrying its
        vertices' min, then an LCT path-min (\(\coderef{c23-summit}\)).
  \item ``connectivity only, edits known offline'' \(\to\) \emph{not} an LCT: rollback DSU /
        segment-tree-on-time (\(\defref{c23-map}\), \(\secref{c23-boundary}\)).
  \item ``path/subtree ranges on a \emph{fixed} tree'' \(\to\) \emph{not} dynamic: HLD (\(\chref{c16}\))
        or Euler+Fenwick (\(\chref{c24}\)) (\(\defref{c23-boundary-pick}\)).
  \item ``counting / aggregating over \emph{all} simple paths of a fixed tree'' \(\to\) centroid
        decomposition (\(\chref{c18}\)), not a per-path fold (\(\secref{c23-boundary}\)).
\end{itemize}

\chaptersection{Mixed Problem Set}

Ten entries, each live-verified against the current Codeforces problem set (names, ratings, and tags
re-pulled this session) and absent from every other roster in the book (the owned IDs -- CF 1601C in
Ch.~19, and this chapter's own SPOJ solveds -- are excluded). None says ``use a link-cut tree'';
several deliberately \emph{do not} want one, and the drill is to say so. Each hint names the
reduction (or the rejection) and the Chapter~25 listing it leans on, ordered ascending by rating.

\begin{mixedset}
  \practice{Chain Queries}{CF 1975E}
    {https://codeforces.com/problemset/problem/1975/E}{CF 2100 L5}
    {A \emph{fixed} tree with queries resolved by a DFS order plus a binary search on the answer --
     no edge is ever linked or cut. Run the top gate of the pick (\(\defref{c23-boundary-pick}\)):
     the topology is static, so this is a static path/chain job for HLD (\(\chref{c16}\)) or an
     Euler tour, \emph{not} a link-cut tree. Use this entry to practise REJECTING the LCT for want
     of a moving edge; contrast \(\secref{c23-boundary}\)'s first spec.}
  \practice{Avoid Rainbow Cycles}{CF 1408E}
    {https://codeforces.com/problemset/problem/1408/E}{CF 2400 L5}
    {Greedy over a graph whose vertices and colours form a bipartite structure, unioning edges with a
     DSU so as to avoid ``rainbow cycles'' -- one maximum-weight forest pass, no online path value.
     The pick hands this to DSU + greedy (\(\chref{c28}\), \(\defref{c23-map}\)): you never query a
     path on a live forest, so an LCT earns nothing. Recognise the Kruskal shape.}
  \practice{Connectivity Addicts}{CF 1738F}
    {https://codeforces.com/problemset/problem/1738/F}{CF 2400 L5}
    {Batched edge insertions followed by connectivity questions -- ``are \(u,v\) in one component?''
     with no path aggregate. That is exactly the connectivity-only row of \(\defref{c23-map}\): a
     plain or rollback DSU (\(\chref{c28}\)) answers each in \(\bigO{\alpha}\), and the tree topology
     is not even being folded. Reject the LCT; the whole cost of \code{makeroot}/\code{access} buys a
     boolean you get free.}
  \practice{Groceries in Meteor Town}{CF 1628E}
    {https://codeforces.com/problemset/problem/1628/E}{CF 3100 L6}
    {Path queries on a \emph{static} tree combined with a binary search / offline ordering -- the
     edges never change. The first question of \(\defref{c23-boundary-pick}\) (does the topology
     change?) is \emph{no}, so route it to a fixed-tree engine (HLD \(\chref{c16}\) or
     Euler+Fenwick \(\chref{c24}\)). The lesson is that ``path query on a tree'' is not by itself an
     LCT trigger -- the \emph{moving} edge is.}
  \practice{Wait for it}{CF 696E}
    {https://codeforces.com/problemset/problem/696/E}{CF 3000 L6}
    {A forest that only \emph{grows} by links while you query aggregates near the frontier -- the
     genuine link-plus-subtree shape of \(\secref{c23-beyond}\). It tempts the \(\mathit{vir}\)
     sibling of \(\coderef{c23-subtree}\); but check whether links are monotone (never cut): if so an
     Euler tour over the \emph{final} tree plus a BIT (\(\chref{c24}\)) is leaner. Decide with
     \(\defref{c23-map}\), then justify the heavier tool if you take it.}
  \practice{Pastoral Oddities}{CF 603E}
    {https://codeforces.com/problemset/problem/603/E}{CF 3000 L6}
    {Edge insertions arrive all known in advance and you maintain a component property, emitting a
     threshold edge weight -- the offline connectivity/divide-and-conquer row of \(\defref{c23-map}\).
     Because you may read every query first, a rollback DSU with D\&C on the timeline
     (\(\chref{c18}\), \(\chref{c28}\)) beats an LCT outright: \(\bigO{\alpha}\) amort per merge and a
     tenth of the code. A textbook ``the LCT is the wrong tool'' drill (\(\secref{c23-boundary}\)).}
  \practice{Surprise me!}{CF 809E}
    {https://codeforces.com/problemset/problem/809/E}{CF 3100 L6}
    {Choose a random subset of marked vertices and aggregate a path quantity over them -- the classic
     \emph{virtual tree} problem: compress the marked set to the smallest subtree spanning it and DP
     on that. This is Chapter~26's territory (\(\chref{c30}\)), \emph{not} link-cut: the base tree is
     fixed and you never rewire it. Forward-reference the virtual-tree construction at
     \(\secref{c23-boundary}\).}
  \practice{Raining season}{CF 1019E}
    {https://codeforces.com/problemset/problem/1019/E}{CF 3200 L6}
    {A fixed tree where each edge is a linear function of a query parameter and you want the
     extreme along a path -- a Li Chao / convex-hull-per-path problem. On a static tree that is HLD
     (\(\chref{c16}\)) with a Li Chao tree at each chain (\(\chref{c29}\)), never an LCT
     (\(\defref{c23-map}\)): the topology is fixed, so the extra evert machinery is dead weight while
     the hard part is the \emph{envelope}, not the path exposure.}
  \practice{Send Tree to Charlie}{CF 1254E}
    {https://codeforces.com/problemset/problem/1254/E}{CF 3300 L6}
    {A per-node state on a \emph{fixed} tree with no route you fold and no edge that moves -- a
     tree-DP / rerooting problem (\(\chref{c10}\), \(\chref{c20}\)). Fourth bullet of
     \(\secref{c23-boundary}\)'s warning box: an LCT needs a path dimension to justify itself; here
     there is none, so halving or evert buys nothing. Recognise ``a state per node, no route'' and
     hand it to DP.}
  \practice{Gregor and the Two Painters}{CF 1548E}
    {https://codeforces.com/problemset/problem/1548/E}{CF 3400 L6}
    {Not a tree problem at all: an \(n\times m\) grid whose cell \((i,j)\) has value
     \(a_i+b_j\) -- a rank-one matrix -- and you binary-search the answer, counting qualifying cells
     with a two-pointer sweep over the sorted rows and columns. It sits in this set purely as the
     sharpest rejection in \(\defref{c23-boundary-pick}\): no graph, no path, no dynamic edge. If
     you cannot immediately say ``nothing here is a forest,'' reread \(\secref{c23-boundary}\).}
\end{mixedset}


\chaptersection{The Summit -- CF 487E Tourists}

\begin{notebox}[Why this is the chapter top rope]
Every in-section problem here ran the LCT on a structure that \emph{already was} a forest. The
Summit hands you a general connected \emph{graph} and asks a path question that is not about a
single route but about \emph{some} simple route -- the exact shape \(\secref{c23-boundary}\)'s last
spec flags: ``condense first, then run the LCT.'' The biconnected (block-cut) tree is the condensation
that makes the question a path fold: once a block node carries the cheapest city inside it, the answer
is a plain path-\(\min\) on a forest -- \code{makeroot} then \code{access} (\(\secref{c23-makeroot}\)).
It is the honest summit, not the highest badge: it is the hardest problem this chapter fully derives
and hand-replays to the judge's output. \emph{One caveat the pick already told you:} the block-cut
tree here is \emph{static} (edges never change, only prices do), so a fixed-tree HLD would also serve;
we instantiate the engine this chapter built -- the LCT -- and treat the reduction, not the dynamism,
as the difficulty (\(\defref{c23-map}\)).
\end{notebox}

\solvedproblem{c23-summit}{Tourists}{CF 487E}
  {https://codeforces.com/problemset/problem/487/E}{CF 3200 L6}
  {Block-cut (biconnected) tree condensation with Tarjan edge-stack; a block node stores the minimum city price among its members and a city node its own price, so the cheapest city on SOME simple a-b route is the path-minimum on the block-cut tree; maintain it online with the chapter's evert-then-access LCT path fold, a price change re-pushing the multiset minima of the (few) blocks containing that city}

\begin{problemstatement}[CF 487E Tourists]
There are \(n\) cities (numbered \(1\) to \(n\)) joined by \(m\) bidirectional roads; city \(i\)
sells a souvenir at price \(w_i\). There are \(q\) online queries of two types:

\textbf{C a w} --- the price in city \(a\) changes to \(w\).\\
\textbf{A a b} --- a tourist travels from city \(a\) to city \(b\) along a \emph{route} (a sequence of
\emph{distinct} cities with consecutive ones joined by a road); he buys the souvenir at the cheapest
city on his route. Report the \emph{minimum possible} price he can achieve, i.e.\
\[
  \min_{\text{simple routes } x_1=a,\dots,x_k=b}\ \min_{1\le i\le k} w_{x_i},
\]
over all routes from \(a\) to \(b\) that visit no city twice (the route may be just \([a]\) when
\(a=b\)).

\textbf{Input.} Line~1: \(n,m,q\) (\(1\le n,m,q\le 10^5\)). Next \(n\) values \(w_i\)
(\(1\le w_i\le 10^9\)). Next \(m\) lines a road \(a_j,b_j\) (\(a_j\ne b_j\), no duplicate road). The
graph is guaranteed connected (a route always exists). Then \(q\) lines, each \texttt{C a w} or
\texttt{A a b}.

\textbf{Output.} For each \texttt{A} query, print the answer on its own line.

\textbf{Sample 1.} triangle \(1\!-\!2\!-\!3\!-\!1\) with prices \(1,2,3\):
\texttt{A 2 3} \(\to\) \texttt{1}; \texttt{C 1 5}; \texttt{A 2 3} \(\to\) \texttt{2}.\\
\textbf{Sample 2.} \(n{=}7\), prices \(1..7\), roads \((1,2),(2,5),(1,5),(2,3),(3,4),(2,4),\)
\((5,6),(6,7),(5,7)\); the four \texttt{A} queries answer \texttt{2, 1, 5, 3}.
\end{problemstatement}

\begin{attemptgate}[Attempt it yourself first -- \lv{6}, budget a real afternoon]
Running a graph search per \texttt{A} query and taking the cheapest city on \emph{one} found path
under-estimates the tourist's freedom: within any biconnected region he can detour. And prices
change, so nothing can be baked once.
\peek{a tourist who enters a 2-vertex-connected region can reach \emph{any} city of that region
without revisiting a city -- what single number therefore summarises a whole region?}
\nudge{Condense the graph to its block-cut tree: a bipartite tree of city-nodes and block-nodes
(maximal biconnected components). Make each block-node's value the \emph{minimum} price among its
member cities and each city-node's value its own price. Then the cheapest achievable on a route
\(a\to b\) is exactly the minimum node value on the unique \(a\)-to-\(b\) path of that tree.}
\fullhint{Build the block-cut tree once with Tarjan's edge stack (block nodes \(n+1,\dots,n+B\));
\(O(n+m)\). Run the chapter's LCT \emph{on that tree} with the path-\(\min\) fold
(\(\mathit{mn}[0]=+\infty\), \(\mathit{mn}\) = min of value and the two children): \texttt{A a b}
\(=\) \code{makeroot(a); access(b); return mn[b]}. Keep, per block, a \code{multiset} of its member
cities' prices so its value is \(*\)set.begin(). A price change \texttt{C a w} updates city-node
\(a\)'s value to \(w\) and, for each of city \(a\)'s \(\bigO{\deg}\) containing block-nodes, erases
the old price from that block's multiset, inserts \(w\), and re-pushes the block-node's value via one
\code{access}. Sum over all block sizes is \(\bigO{n+m}\), so an average update touches few nodes;
the total is \(\bigO{(n+m)+q\log n\cdot\bar{d}}\). Read/write 64-bit prices. \emph{Notice the
honesty check:} the tree itself never changes, so a fixed-tree HLD + segment tree is an equally valid
-- and by \(\defref{c23-boundary-pick}\) the leaner -- choice; the reduction to a block-cut tree is
the actual hard part.}
\end{attemptgate}

\solstep{Understanding}
A ``route'' is a \emph{simple path}, and the tourist wants the cheapest city he can arrange to pass
through. The pivotal structural fact is that a biconnected component (block) is \emph{enterable and
freely movable}: for any three distinct cities \(x,y,z\) in a 2-connected graph there is a simple
path from \(x\) to \(y\) through \(z\). So a tourist who must traverse a block from entry city
\(p\) to exit city \(q\) can also swing by \emph{any} city of that block -- in particular its
cheapest one -- without ever repeating a city. Hence a block contributes the single number
\(\min\{w_v : v\in\text{block}\}\), and only its articulation cities constrain the route.

Form the block-cut tree: bipartite, city-nodes for the original cities, block-nodes for the blocks,
a city-node joined to a block-node when the city belongs to the block. Every simple route \(a\to b\)
in the graph projects to the unique \(a\)-to-\(b\) path in the block-cut tree (the articulation
cities and blocks it must cross), and conversely the block-cut path is realisable as a graph route.
The cheapest city obtainable is therefore the minimum node value along that tree path, where a city
node is valued by \(w\) and a block node by its members' minimum. This is a path-\(\min\) on a tree
with point updates -- precisely \(\secref{c23-makeroot}\)'s fold.

\solstep{Key Observations}
\begin{observation}
\textbf{Block mobility.} Inside one biconnected component, entering at \(p\) and leaving at \(q\),
the tourist can pass through \emph{any chosen city} of the component on a simple path. Thus the
component is summarised by \emph{one} value -- the minimum member price -- and no route ever needs
to know more about it. (This collapses a whole block to a single node weight.)
\end{observation}
\begin{observation}
\textbf{Answer \(=\) path-min on the block-cut tree.} A route from \(a\) to \(b\) corresponds to the
unique \(a\)-to-\(b\) path in the block-cut tree; the minimum achievable price equals
\(\min\{\text{node value} : \text{node on that tree path}\}\) -- city nodes give their own price,
block nodes give their members' minimum. The path fold \(\mathit{mn}\) of \(\thmref{c23-path-query}\)
answers it in \(\bigO{\log n}\).
\end{observation}
\begin{observation}
\textbf{A price change touches few nodes.} Changing city \(a\)'s price to \(w\) updates city-node
\(a\), and each block-node containing \(a\) whose multiset-minimum may move -- that is
\(\bigO{\deg_{\text{BCT}}(a)}\) block nodes, each maintained by a \code{multiset} so the new value is
\(*\)begin. The block-cut \emph{topology} is fixed; only node values change, which is exactly what an
LCT point-update (access, set, pull) handles.
\end{observation}

\solstep{Think It With Me}
\begin{thinkbox}
\deadend{For each \texttt{A a b}, run BFS/Dijkstra and take the cheapest city on the found shortest path}
{Two failures. First, cost is not the number of roads, so a shortest path is the wrong object.
Second, and fatal, the tourist may take a \emph{longer} simple route to touch a cheap city a shortest
path skips. The answer is not ``min over one canonical path'' but ``min over the best simple route,''
which the block structure makes a tree-path problem. \textbf{WA}.}
\deadend{Value a block-node by only its articulation cities}
{A block's cheapest city need not be an articulation point -- it can be a leaf interior to the block,
which the tourist can still reach (Observation~1). Valuing the block by just its boundary cities
misses the whole interior and gives \textbf{WA}; the node value must be the minimum over \emph{every}
member city.}
\deadend{Rebuild the block-cut tree or DFS the tree path on every query}
{The tree has \(\bigO{n+m}\) nodes and there are \(q\le 10^5\) queries -- a per-query walk is
\(\bigO{q(n+m)}\), \tle. You need the path-\(\min\) to be polylogarithmic and the price updates to be
local, which is the LCT's job; but note the topology never changes, so HLD would also be correct
(\(\defref{c23-map}\)) -- the dynamism is in the \emph{values}, not the edges.}
\aha{Condense \emph{once} to the block-cut tree (Tarjan edge stack), give each block-node its
members' multiset-minimum and each city-node its own price, then answer \texttt{A} with the chapter's
\code{makeroot(a); access(b); mn[b]} path-\(\min\), and \texttt{C} by re-pushing the \(\bigO{\deg}\)
affected block-nodes. The hard insight is the \emph{reduction}; the data structure is the fold you
already built in \(\secref{c23-makeroot}\).}
\end{thinkbox}

\solstep{The Approach Ladder}
\begin{center}\renewcommand{\arraystretch}{1.25}
\begin{tabular}{@{}c p{0.23\linewidth} p{0.29\linewidth} l c@{}}
\toprule
\textbf{\#} & \textbf{Approach} & \textbf{Idea} & \textbf{Time} & \textbf{Verdict}\\
\midrule
1 & \lv{2}\ BFS the graph per query & min city on one shortest path & \(\bigO{q(n+m)}\) & \tle \\
2 & \lv{3}\ BFS but allow any simple path & interiors are reachable, missed & -- & \textcolor{warningred}{\ding{55}\ \textbf{WA}} \\
3 & \lv{4}\ block-cut tree, DFS path-min each query & right model, linear query & \(\bigO{q(n+m)}\) & \tle \\
4 & \lv{6}\ \textbf{block-cut tree + LCT path-min} & condense once, then fold & \(\bigO{q\log n}\) & \ac \\
\bottomrule
\end{tabular}\end{center}
Rung~4 is the Summit: rung~3 has the right \emph{model} but a linear query; replacing the per-query
tree walk with the LCT path fold (\(\thmref{c23-path-query}\)) is the whole win, and the price update
becomes a local multiset re-push.

\solstep{Solution}
Read the graph and prices. Run Tarjan's biconnected-component DFS with an edge stack: whenever
\(\mathit{low}[v]\ge \mathit{dfn}[u]\), pop the stack down to the current tree edge -- those edges'
distinct endpoints are one block, given a fresh node id \(n+\mathit{bcnt}\). Link each member city to
the block node (in the LCT as well as the tree), insert its price into the block's multiset, and set
the block node's value to the multiset minimum; set each city node's value to its own price. Answer
\texttt{A a b} by \code{makeroot(a); access(b)} and reading \(\mathit{mn}[b]\). On \texttt{C a w}, for
each block containing \(a\) erase the old price and insert \(w\), then re-push that block node's value;
finally re-push city node \(a\)'s value. Everything is a path-\(\min\) with point updates, exactly the
\(\secref{c23-makeroot}\) engine.

\solstep{Proof of Correctness}
\begin{proofbox}
\textbf{(Route \(\Rightarrow\) tree path, upper bound.)} Let \(\mathcal{P}\) be the unique
\(a\)-to-\(b\) path in the block-cut tree. Any simple graph route \(R\) from \(a\) to \(b\) crosses
exactly the blocks and articulation cities of \(\mathcal{P}\) (removing an articulation city
disconnects the graph along \(\mathcal{P}\)), so every city of \(R\) lies in some block-node or is a
city-node of \(\mathcal{P}\). Hence \(\min_{v\in R} w_v \ge \min_{x\in\mathcal{P}} \mathit{val}(x)\):
a city is either a city-node (its value is \(w\)) or inside a block whose node value is the block's
\(\min\le w\). So no route beats \(\min\mathcal{P}\). \textbf{(Tree path \(\Rightarrow\) route, lower
bound.)} Write \(\mathcal{P}=a, B_1, c_1, B_2,\dots,b\). Concatenate, for each block \(B_i\), a simple
path inside \(B_i\) from its entry articulation to its exit articulation that \emph{passes through}
\(B_i\)'s cheapest city -- realisable because a block is 2-connected (Observation~1), and these sub-paths
meet only at articulation cities, so their union is a genuine simple route. Choosing the cheapest city
of the block attaining \(\min\mathcal{P}\) shows the value \(\min\mathcal{P}\) is achievable. \textbf{(Updates.)}
The block-cut tree is topologically fixed (the graph's edges never change), so \(\mathcal{P}\) is a
function of only \(a,b\); a price change alters one city-node value and, per containing block, that
block's multiset minimum, which is exactly the set of node values on future paths that the change can
affect. The LCT \(\mathit{mn}\) fold returns \(\min\mathcal{P}\) after each update, hence every
\texttt{A} prints the true minimum. \(\Box\)
\end{proofbox}

\solstep{Implementation}
\begin{lstlisting}[caption={CF 487E Summit: Tarjan block-cut tree (edge stack) built once, then the chapter's evert-then-access LCT path-min over it. A block node stores the multiset-min of its member city prices, a city node its own price; A a b = makeroot(a);access(b);mn[b], C a w re-pushes the few blocks containing a plus city a itself. ASCII only, identifiers without underscores.}, label={lst:c23-summit}]
#include <bits/stdc++.h>
using namespace std;
typedef long long ll;
const int MAXN = 200005;                 // n cities + <= n blocks
const ll  INFL = (ll)4e18;
int n, m, q;
// ---------- LCT with path-MIN ----------
int ch[MAXN][2], fa[MAXN]; bool revTag[MAXN];
ll val[MAXN], mn[MAXN];
bool isRoot(int x){ int f=fa[x]; return f==0 || (ch[f][0]!=x && ch[f][1]!=x); }
void pull(int x){ mn[x]=val[x];
    if(ch[x][0]) mn[x]=std::min(mn[x],mn[ch[x][0]]);
    if(ch[x][1]) mn[x]=std::min(mn[x],mn[ch[x][1]]); }
void pushRev(int x){ if(!x) return; std::swap(ch[x][0],ch[x][1]); revTag[x]^=1; }
void push(int x){ if(revTag[x]){ pushRev(ch[x][0]); pushRev(ch[x][1]); revTag[x]=0; } }
void pushPath(int x){ if(!isRoot(x)) pushPath(fa[x]); push(x); }
void rot(int x){ int y=fa[x], z=fa[y], d=(ch[y][1]==x);
    if(!isRoot(y)) ch[z][ch[z][1]==y]=x; fa[x]=z;
    ch[y][d]=ch[x][d^1]; if(ch[x][d^1]) fa[ch[x][d^1]]=y;
    ch[x][d^1]=y; fa[y]=x; pull(y); pull(x); }
void splay(int x){ pushPath(x);
    while(!isRoot(x)){ int y=fa[x], z=fa[y];
        if(!isRoot(y)) rot( (ch[y][0]==x)==(ch[z][0]==y) ? y : x );
        rot(x); } }
void access(int x){ for(int last=0; x; last=x, x=fa[x]){ splay(x); ch[x][1]=last; pull(x);} }
void makeroot(int x){ access(x); pushRev(x); }
void link(int u,int v){ makeroot(u); fa[u]=v; }
void setVal(int x,ll w){ access(x); val[x]=w; pull(x); }
ll  pathMin(int u,int v){ makeroot(u); access(v); return mn[v]; }
// ---------- Block-cut tree (Tarjan edge stack) ----------
vector<pair<int,int> > g[MAXN]; int eu[MAXN], ev[MAXN];
int dfn[MAXN], low[MAXN], tmr, estk[MAXN], etop, bcnt;
ll wcity[MAXN];
multiset<ll> ms[MAXN];                     // per block node: member city prices
vector<int> ofBlock[MAXN];                 // city -> block nodes containing it
void dfs(int u,int pe){ dfn[u]=low[u]=++tmr;
    for(size_t k=0;k<g[u].size();k++){ int v=g[u][k].first, id=g[u][k].second;
        if(!dfn[v]){ estk[++etop]=id; dfs(v,id); low[u]=std::min(low[u],low[v]);
            if(low[v]>=dfn[u]){                       // peel one block
                ++bcnt; int node=n+bcnt; vector<int> verts;
                while(true){ int e=estk[etop--];
                    if(eu[e]!=0){verts.push_back(eu[e]);} // dedup below
                    verts.push_back(ev[e]);
                    if(e==id) break; }
                sort(verts.begin(),verts.end()); verts.erase(unique(verts.begin(),verts.end()),verts.end());
                for(size_t t=0;t<verts.size();t++){ int v=verts[t];
                    ms[node].insert(wcity[v]); ofBlock[v].push_back(node); link(v,node); }
                val[node]=*ms[node].begin(); pull(node);
            }
        } else if(id!=pe && dfn[v]<dfn[u]){ estk[++etop]=id; low[u]=std::min(low[u],dfn[v]); }
    } }
int main(){ scanf("%d%d%d",&n,&m,&q);
    for(int i=1;i<=n;i++){ scanf("%lld",&wcity[i]); val[i]=wcity[i]; mn[i]=wcity[i]; }
    for(int i=1;i<=m;i++){ int a,b; scanf("%d%d",&a,&b); eu[i]=a; ev[i]=b;
        g[a].push_back(make_pair(b,i)); g[b].push_back(make_pair(a,i)); }
    for(int i=1;i<=n;i++) if(!dfn[i]) dfs(i,-1);      // build BCT + LCT forest once
    for(int qi=0;qi<q;qi++){ char t; int a,b; ll w;
        scanf(" %c%d%d",&t,&a,&b);
        if(t=='A'){ printf("%lld\n", pathMin(a,b)); }
        else { w=(ll)b; ll old=wcity[a];
            for(size_t k=0;k<ofBlock[a].size();k++){ int node=ofBlock[a][k];
                multiset<ll>::iterator it=ms[node].find(old); if(it!=ms[node].end()) ms[node].erase(it);
                ms[node].insert(w); setVal(node,*ms[node].begin()); }
            wcity[a]=w; setVal(a,w); } }
    return 0; }
// The graph is condensed ONCE (topology fixed); A is a makeroot+access path-min, C is a local
// multiset re-push on the few blocks touching the city. mn[0] unused (pull reads only real children).
\end{lstlisting}

\solstep{Trace on the Official Sample}
\begin{tracebox}[Sample 2: condense, then replay the four A queries as path-minima]
Prices \(w_i=i\). The three triangles are blocks; cities \(2\) and \(5\) are the articulations.
Block nodes: \(B_a{=}8=\{1,2,5\}\) min \(=1\); \(B_b{=}9=\{2,3,4\}\) min \(=2\);
\(B_c{=}10=\{5,6,7\}\) min \(=5\). Tree edges: \(8\!-\!1,8\!-\!2,8\!-\!5,9\!-\!2,9\!-\!3,9\!-\!4,
10\!-\!5,10\!-\!6,10\!-\!7\). Node values: city \(i\mapsto i\); \(8\mapsto1,9\mapsto2,10\mapsto5\).

\textbf{A 2 3.} Tree path \(2\!-\!9\!-\!3\); values \(2,2,3\); min \(=\boxed{2}\). \(\checkmark\)

\textbf{A 6 4.} Path \(6\!-\!10\!-\!5\!-\!8\!-\!2\!-\!9\!-\!4\); values \(6,5,5,1,2,2,4\);
min \(=\boxed{1}\) (realised by detouring through city \(1\) inside block \(8\)). \(\checkmark\)

\textbf{A 6 7.} Path \(6\!-\!10\!-\!7\); values \(6,5,7\); min \(=\boxed{5}\). \(\checkmark\)

\textbf{A 3 3.} Path is the single node \(3\); value \(w_3=\boxed{3}\). \(\checkmark\) -- matching the
judge exactly. (Sample 1: the whole triangle is one block \(B\) with min \(=\min(1,2,3)=1\);
\texttt{A 2 3} walks \(2\!-\!B\!-\!3\) giving \(\min(2,1,3)=1\); after \texttt{C 1 5} the block min
becomes \(\min(5,2,3)=2\), so \texttt{A 2 3} gives \(\min(2,2,3)=2\).)
\end{tracebox}

\solstep{Complexity}
\begin{complexitybox}
Tarjan block-cut construction is \(\bigO{n+m}\) time and space; the tree has
\(\le n+m\le 2\cdot10^5\) nodes (each block costs one node, total block sizes \(\bigO{n+m}\)). Each
\texttt{A} is one \code{makeroot}/\code{access}, \(\bigO{\log N}\) amortised (\(\thmref{c23-splay-amort}\)).
Each \texttt{C a w} touches city node \(a\) plus its \(\bigO{\deg_{\text{BCT}}(a)}\) block nodes, each a
\(\bigO{\log}\) multiset edit and a \(\bigO{\log N}\) LCT \code{setVal}; summing over all price changes
that is \(\bigO{\sum_a \deg(a)}=\bigO{n+m}\) per full sweep, so on typical data the queries cost
\(\bigO{q\log N}\) and the worst single city (a hub in many blocks) is \(\bigO{\deg\log N}\) --
comfortably inside 2 seconds at \(10^5\). Prices up to \(10^9\) fit in \code{long long};
\(\mathit{mn}\) is a genuine \(\min\), and the null child is read only after checking it exists, so no
neutral-value bug can surface.
\end{complexitybox}

\begin{hardeningbox}
\begin{itemize}[nosep,leftmargin=1.4em]
  \item \textbf{A block's value is the min over \emph{all} its cities, not just articulations.} A
        leaf interior to a block is reachable (2-connectivity), so restricting the min to boundary
        cities is the classic \textbf{WA} on this problem (\(\probref{c23-summit}\)'s second dead end).
  \item \textbf{Re-push the containing blocks, then the city, on \texttt{C}.} The block multiset min
        \emph{and} the city node both changed; forgetting the city node misses routes that stop at
        \(a\) itself, forgetting a block misses the interior-min detour (\(\coderef{c23-summit}\)).
  \item \textbf{Pop the edge stack \emph{down to the current tree edge} and dedup endpoints.} A block
        is a set of edges; its vertex set is their distinct endpoints. Missing the \(\mathit{low}\ge
        \mathit{dfn}\) boundary or forgetting to dedup double-counts a city into one block's multiset.
  \item \textbf{Keep \(\mathit{mn}\) a pure \(\min\) and guard null children.} Reading
        \(\mathit{mn}[0]\) as \(0\) (not \(+\infty\)) would sink every path-min to \(0\); the fold is
        a min, so the neutral is \(+\infty\) and \code{pull} must only fold existing children.
  \item \textbf{The topology is static -- but that does not excuse a wrong engine claim.} Because no
        edge is ever added, an honest reader notes HLD + a segment tree also passes
        (\(\defref{c23-boundary-pick}\)); the LCT is used to reuse \(\secref{c23-makeroot}\), not
        because the forest moves. Say so under pressure rather than over-claiming dynamism.
\end{itemize}
\end{hardeningbox}

\begin{recallbox}[Chapter 25 in ten seconds]
A splay shuffles the touched key to the root in \(\bigO{\log n}\) amortised (\(\thmref{c23-splay-amort}\));
a link-cut tree chains splays along preferred paths so \code{access} exposes a root-to-node path in
one loop (\(\defref{c23-access}\), \(\coderef{c23-lct-core}\)). Rooted, the two-access climb reads an
LCA (\(\thmref{c23-access-lca}\), CF DYNALCA \(\coderef{c23-dynalca}\)); with \code{makeroot} and a
reverse tag, any path fold drops out of \code{access} (\(\thmref{c23-path-query}\),
\(\coderef{c23-evert}\)); a virtual-child accumulator extends it to subtree aggregates
(\(\thmref{c23-vsum}\), \(\coderef{c23-subtree}\)). Fixed trees want HLD/Euler, offline connectivity
wants DSU, all-pairs counting wants centroid (\(\defref{c23-boundary-pick}\)). Summit CF 487E
(\(\probref{c23-summit}\)) condenses a graph to its block-cut tree and reads a path-min.
\selftest{
  \item Why is \(\mathit{isRoot}(x)\) not simply \(\mathit{fa}[x]=0\)? \hfill\lv{4}
  \item State the two-access LCA procedure and why the last splayed node is the ancestor. \hfill\lv{5}
  \item Why must \code{pushPath} run before the zig-zig loop in \code{splay}? \hfill\lv{5}
  \item What does \(\mathit{vir}[x]\) hold, and why does the path fold omit it while the subtree fold
        includes it? \hfill\lv{5}
  \item In the Summit, what value does a block node carry and why is it not just its articulations? \hfill\lv{6}
}
\end{recallbox}


\chaptersection{Tool Kit Inventory}

The machine assembled through Chapter~25. Earlier chapters' rows are carried forward; this
chapter's are appended. Pointers only.

\begin{center}
\small
\renewcommand{\arraystretch}{1.25}
\begin{tabular}{@{}p{0.23\linewidth} l c p{0.22\linewidth} >{\raggedright\arraybackslash\scriptsize}p{0.18\linewidth}@{}}
\toprule
\textbf{Tool} & \textbf{Since} & \textbf{Lv} & \textbf{Trigger phrase} & \textbf{Where} \\
\midrule
\multicolumn{5}{l}{\itshape Carried forward} \\
Euler tour / subtree-as-interval & 5 & \lv{3} & ``a fixed tree's subtree is one contiguous range'' & \chref{c24} \\
LCA + binary lifting & 16 & \lv{4} & ``the static-tree ancestor the rooted LCA replaces'' & \chref{c11} \\
Heavy-light decomposition & 22 & \lv{5} & ``a path on a tree that never changes'' & \chref{c16} \\
DSU / rollback DSU & 17 & \lv{4} & ``connectivity only, offline'' & \chref{c28} \\
Centroid decomposition & 24 & \lv{5} & ``counting over ALL paths of a fixed tree'' & \chref{c18} \\
\midrule
\multicolumn{5}{l}{\itshape Chapter 25 -- new this chapter} \\
Splay tree (amortised) & 25.1 & \lv{4} & ``shuffle the touched key up; order statistics'' & \defref{c23-splay}, \coderef{c23-splay-rot} \\
Top-down splay & 25.1 & \lv{4} & ``one downward pass, no parent pointers'' & \coderef{c23-topdown} \\
Preferred paths + access & 25.2 & \lv{4} & ``expose root-to-x on a moving forest in one loop'' & \defref{c23-access}, \coderef{c23-lct-core} \\
Two-access LCA & 25.2 & \lv{5} & ``LCA / connectivity on a growing rooted forest'' & \thmref{c23-access-lca}, \coderef{c23-dynalca} \\
Evert + path fold & 25.3 & \lv{5} & ``sum/min over the current u--v path, rerootable'' & \thmref{c23-path-query}, \coderef{c23-evert} \\
Virtual-child subtree fold & 25.4 & \lv{5} & ``a subtree aggregate on a moving tree'' & \thmref{c23-vsum}, \coderef{c23-subtree} \\
The dynamic-tree engine pick & 25.5 & \lv{5} & ``is the topology even moving, and do I need a path value?'' & \defref{c23-boundary-pick}, \defref{c23-map} \\
Link-cut kit (one listing) & 25 & \lv{5} & ``all the skeletons in one place'' & \coderef{c23-kit} \\
Summit: block-cut + LCT path-min & Summit & \lv{6} & ``cheapest city on some simple route, prices change'' & \coderef{c23-summit} \\
\bottomrule
\end{tabular}
\end{center}

\chaptersection{Related Topics}

\begin{itemize}[nosep, leftmargin=1.4em]
  \item \chref{c16} (Ch.~22, Heavy-Light Decomposition) --- contrast and promise \textbf{discharged}:
        HLD folds a path on a \emph{fixed} tree in \(\bigO{\log^2 n}\); \(\secref{c23-makeroot}\) folds a
        path on a \emph{moving} forest in \(\bigO{\log n}\) amortised. \(\secref{c16-boundary}\) handed
        the moving-topology case forward to this chapter; \(\defref{c23-boundary-pick}\) makes the split
        explicit -- the tell is whether an edge ever appears or vanishes.
  \item \chref{c11} (Ch.~16, LCA) --- backward/contrast: binary lifting answers a static LCA in
        \(\bigO{\log n}\) but cannot survive a link or cut; the two-access LCA
        (\(\thmref{c23-access-lca}\), \(\coderef{c23-dynalca}\)) is its dynamic sibling, recomputing the
        ancestor as the rooted forest rewires.
  \item \chref{c18} (Ch.~24, Centroid Decomposition) --- contrast: centroid counts over \emph{all} paths
        of a fixed tree; the LCT folds \emph{one} named path of a changing tree. Both die if you mix the
        scopes -- centroid needs a static topology, LCT needs a moving one (\(\defref{c23-map}\)).
        \textbf{This is the Previous chapter.}
  \item \chref{c24} (Ch.~5, Euler Tour / tree representations) --- backward: the Euler tour is the
        \emph{static} way a subtree becomes an interval; an Euler-tour tree keeps that tour in a balanced
        BST and is the simpler dynamic-forest tool when you only need connectivity, not a path fold
        (\(\secref{c23-beyond}\)).
  \item \chref{c28} (Ch.~17, DSU / Sack) --- forward/contrast: a rollback DSU answers offline
        connectivity in \(\bigO{\alpha}\) with a fraction of the code; reach for it before an LCT whenever
        the edits are all known in advance (\(\defref{c23-map}\)).
  \item \chref{c29} (Ch.~23, Li Chao / Segment-Tree Beats) --- forward/contrast: a path whose aggregate is
        a line envelope evaluated at a variable is Li Chao hosted on HLD chains (CF 1019E in the mixed
        set), not a monoid fold on an LCT.
  \item \chref{c10} (Ch.~13, Tree DP) --- contrast: a state \emph{per node} with no route dimension is DP
        / rerooting (\(\chref{c20}\)); the LCT needs a path (or subtree) you fold, so halving/evert buys
        nothing on pure DP specs (\(\secref{c23-boundary}\)).
  \item \chref{c30} (Ch.~26, Virtual Trees) --- forward: the mixed-set CF 809E hands a random-marked-subset
        aggregate to the virtual tree, the next chapter's engine -- compressing a fixed tree to the nodes
        that matter rather than making it dynamic. \textbf{This is the Next chapter.}
  \item \chref{c21} (Ch.~27, Mo's Algorithm) --- forward: an offline batch of node-set queries with no
        inverse is Mo-on-trees, complementing the LCT's online, invertible path folds
        (\(\secref{c23-boundary}\)).
\end{itemize}

\begin{lastword}[Last Word --- Chapter 25]
\emph{``A splay tree remembers what you touched last and charges you for it in the right order; a
link-cut tree chains those splays along preferred paths and re-preferences them with a single access,
so a forest that keeps rewiring its edges can still answer a path -- and, with a virtual-child ledger,
a subtree -- in amortised \(\bigO{\log n}\).''}

\medskip
Chapter~22 folded paths only on a tree that stood still; this chapter dropped that assumption. Section
25.1 built the splay and its potential-function bound (\(\thmref{c23-splay-amort}\)); 25.2 turned
splays into a forest with \code{access} and read an LCA two-accesses at a time
(\(\thmref{c23-access-lca}\)); 25.3 added \code{makeroot} and a lazy reversal so any path fold drops out
of \code{access} (\(\thmref{c23-path-query}\)); 25.4 paid the honest price of a \emph{subtree} on a
moving tree with the \(\mathit{vir}\) ledger (\(\thmref{c23-vsum}\)); and the Summit CF 487E
(\(\probref{c23-summit}\)) showed that sometimes the tree you need is one you must \emph{build} --
condensing a graph to its block-cut tree before the fold even applies.

\medskip
Carry three reflexes forward:
\begin{enumerate}[nosep, leftmargin=1.6em]
  \item \emph{Ask whether the topology actually moves.} If no edge is ever linked or cut, you wanted
        HLD or Euler+Fenwick, not an LCT (\(\defref{c23-boundary-pick}\)); the moving edge is the whole
        licence to pay splay constants.
  \item \emph{Access is the one primitive.} \code{findroot}, \code{makeroot}, \code{link}, \code{cut}
        and every path or subtree fold are the same root-to-node exposure wearing different hats
        (\(\defref{c23-access}\), \(\coderef{c23-kit}\)).
  \item \emph{Push your lazy tags before you rotate.} A splay that rotates under an un-pushed reverse
        tag reattaches mirrored children -- the deepest, rarest WA in the structure
        (\(\defref{c23-makeroot}\), \(\coderef{c23-evert}\)).
\end{enumerate}

\medskip
And three habits, drilled until boring: keep \(\mathit{isRoot}\) a splay test, never \code{fa==0}; keep
the path fold \(\mathit{sum}\) and the subtree fold \(\mathit{sub}\) as disciplined siblings so a query
can never grab the wrong one; and when a ``dynamic tree'' problem hands you a \emph{graph}, reach for a
condensation (block-cut or spanning forest) before you reach for \code{makeroot}. You now own the one
tree engine that answers while the tree itself changes -- the dynamic companion to every static
structure in Part~V.

\medskip
\textbf{Next:} \chref{c30} (Virtual Trees) turns the difficulty inside-out. Instead of a
\emph{moving} tree with all of \(n\) nodes, you get a \emph{fixed} tree and a small marked subset per
query, and the win is to compress the tree down to just the marked nodes and their LCAs -- an
\(\bigO{k\log k}\) skeleton -- then run a DP on it. Where this chapter made the structure move, the
next one makes it \emph{shrink}, and it cashes in the LCA machinery of Chapter~16 that this chapter
merely borrowed.
\end{lastword}


\closechapter
